% ============================================================
% Erster Codierdurchgang (Phase 3 nach Kuckartz, 2018)
% Interview: IP-01 | Erstellt: 20.06.2026 | Modell: claude-sonnet-4-6
% Sinneinheiten: 41 | Kodierungen: 69
% HK0:9 HK1:24 HK2:24 HK3:12
% HINWEIS: Diese Datei ist nicht git-versioniert (Datenschutz).
% Alle Kodierungen sind KI-generierte Arbeitshypothesen
% und muessen vor Verwendung manuell geprueft werden.
% ============================================================

\section*{Interview \texttt{IP-01} -- Erster Codierdurchgang}

\noindent\textit{Erstellt am 20.06.2026 | Sinneinheiten: 41 |
Kodierungen: 69 (HK0: 9, HK1: 24,
HK2: 24, HK3: 12)}\\[0.5em]

\noindent\textbf{Hinweis:} Diese Kodierungen wurden KI-gest\"utzt erstellt und
sind als Arbeitshypothesen f\"ur den ersten Codierdurchgang zu verstehen.
Alle Zuweisungen m\"ussen manuell gepr\"uft und ggf. korrigiert werden
\parencite{kuckartzQualitativeInhaltsanalyseMethoden2018}.

% ── HK0: Übergreifende Kontextfaktoren ──────────────────────────────────────────────
\subsection*{HK0 -- Übergreifende Kontextfaktoren}

\begin{description}
  \item[\texttt{\#}\textbf{00:02:18}\texttt{\#} (Sinneinheit 2)]
  \begin{quote}
  Ja, wir haben bei uns im Haus vor Jahren uns dazu entschieden, eine eigene Abteilung zu machen mit dem KP-01. Diese Abteilung ist so das Herzstück für unsere Digitalisierung in Form von KI, Big Data und alle Möglichkeiten, Robotik, Automatisierung auch. Und haben wirklich damals gesagt, es ist besser, bevor wir dieses Wissen gleich versuchen, breit in der Organisation zu verankern, dass man mal quasi das Herz all dieser Vorhaben zum Schlagen bringt. Und der KP-01 hat den Auftrag, mit all seinen Themen, die er strategisch angeht, uns in weiterer Folge zu befähigen. Und wir haben einfach gesagt, dieses Herz soll in der Mitte sein. Und rundherum sollen wir in unseren Organisationseinheiten dann das Know-how aufbauen und dieses Wissen, was dort primär generiert wird, dann in die breite Organisation tragen.
  \end{quote}
  \textbf{Kategorie:} HK0 -- Übergreifende Kontextfaktoren\\
  \textbf{Begr\"undung:} Die Passage beschreibt die organisationale Struktur und den institutionellen Einführungsprozess von KI im Unternehmen, einschließlich der strategischen Entscheidung zur Gründung einer zentralen KI-Abteilung. Dies ist ein relevanter organisationaler Rahmenbedingungsfaktor, der für das Verständnis der motivationalen Dynamik der Führungskraft bedeutsam ist, ohne direkt auf ein spezifisches Grundbedürfnis zu verweisen.
  \\\textit{Memo: Die Passage ist rein kontextuell-beschreibend und schildert den institutionellen KI-Einführungsprozess (zentrale Einheit als 'Herzstück', Befähigungsauftrag, Wissenstransfer in die Organisation). Es ist kein direkter Bezug zu Autonomie-, Kompetenz- oder Eingebundenheitserleben der Interviewperson erkennbar. HK0 ist angemessen, da die organisationale Rahmenbedingung für spätere motivationale Aussagen als Hintergrundfolie dient.}
\end{description}
\medskip

\begin{description}
  \item[\texttt{\#}\textbf{00:02:18}\texttt{\#} (Sinneinheit 3)]
  \begin{quote}
  Es hat sich dann aber auch so entwickelt in weiterer Folge, dass das, was der KP-01 macht, sehr, sehr gut ist. Aber wir wollen jetzt so viele eigeninteressierte Mitarbeiterinnen und Mitarbeiter in unseren Bereichen haben, die sich auch dann selber Wissen aneignen. Wir haben dazu dann Schulungen auch aufgesetzt. M365, Co-Pilot, beiden Schulungen hat es auch bei uns im Haus gegeben. Einfach wirklich Fokus-Schulung zu Bau-API, aber immer mit KI-Kontext. Und wenn ich da jetzt sagen würde, Regeln gibt es logischerweise bitte auch, sowohl Richtung IT-Security und alles, was mit KI zu tun hat, ist bei uns geregelt im Haus. In der Praxis ist es dann so aber, dass viele Leute einfach Learning by Doing selber Sachen ausprobieren wollen. Und somit kann ich die Frage beantworten. Wir haben Regeln, wir haben Schulungen, wir haben es organisatorisch gelöst.
  \end{quote}
  \textbf{Kategorie:} HK0 -- Übergreifende Kontextfaktoren\\
  \textbf{Weitere Kategorien:} HK1\\
  \textbf{Begr\"undung:} Die Passage beschreibt den institutionellen KI-Einführungsprozess der Organisation, einschließlich organisatorischer Rahmenbedingungen wie Schulungsangebote (M365, Co-Pilot), IT-Security-Regeln und struktureller Lösungen. Diese Informationen sind kontextuell relevant für das Verständnis der motivationalen Dynamik, beziehen sich aber primär auf organisationale Rahmenbedingungen ohne direkte Zuordnung zu einem der drei Grundbedürfnisse.
  \\\textit{Memo: Die Passage ist primär deskriptiv-kontextuell gehalten und schildert die organisationale KI-Einführungsarchitektur. Der Autonomiebezug ist indirekt, da er sich auf Mitarbeitende und nicht explizit auf die Führungskraft selbst bezieht. Für spätere Analysephasen relevant: Spannung zwischen institutionellen Regeln und dem Wunsch nach selbstgesteuertem Lernen (Autonomie vs. Kontrolle).}
\end{description}
\medskip

\begin{description}
  \item[\texttt{\#}\textbf{00:04:52}\texttt{\#} (Sinneinheit 5)]
  \begin{quote}
  In meiner Funktion habe ich vorwiegend, also mein Hauptziel ist es immer, dass ich meine sechs Abteilungsleiter bestmöglich unterstütze und ihnen Rahmenbedingungen für Spitzenerfolge ermögliche, um es sehr hochtrabend zu sagen. Sprich, ich schaue, dass alle meine sechs Abteilungsleiter momentan nur Herren noch einfach bestmöglich ihre eigene Aufgabe und Verantwortung erfüllen können, dass sie dazu kompetent genug sind und dass, wenn es Probleme gibt, die ich für sie am besten regeln kann, dass sie dann sofort zu mir kommen, dass ich diese Probleme regle. Dadurch habe ich viele Besprechungen. Ich selber habe den Luxus, dass meine Leute so viel ohne mich machen, dass ich circa 50 Prozent meiner Zeit wirklich für persönliche Themen in Form von beruflichen Themen widmen kann, Mitarbeit in Projekten auf Bundes-, Landes-, aber auch Unternehmensebene hier in der Bank.
  \end{quote}
  \textbf{Kategorie:} HK0 -- Übergreifende Kontextfaktoren\\
  \textbf{Weitere Kategorien:} HK1\\
  \textbf{Begr\"undung:} Die Passage beschreibt die organisationale Rolle und Führungsstruktur der Interviewperson (sechs Abteilungsleiter, Projekttätigkeit auf verschiedenen Ebenen) und liefert damit relevanten biografisch-institutionellen Kontext für das Verständnis der späteren motivationalen Dynamik beim KI-Einsatz.
  \\\textit{Memo: Die Sinneinheit enthält keinen direkten Bezug zu generativer KI, ist jedoch als Kontextpassage motivational relevant, da sie den Ausgangszustand der Autonomie und Führungsverantwortung beschreibt, vor dem spätere KI-bezogene Erfahrungen interpretiert werden müssen.}
\end{description}
\medskip

\begin{description}
  \item[\texttt{\#}\textbf{00:10:14}\texttt{\#} (Sinneinheit 11)]
  \begin{quote}
  Wir hatten ja eine große Ausschreibung an externe Partner und hatten dann neun Angebote vorliegen. Ich bin von meinem Urlaub an dem Tag zurückgekommen, wo die Deadline war, für diese neun sehr hochwertigen Angebote. Es waren über 566 Folien, die ich da von diesen neun Anbietern bekommen habe. Und ich habe da im Vorfeld, weil ich ja gewusst habe, das ist erstens einmal zeitlich ein sehr, sehr knappes Thema. Freitag war die Deadline. Montag musste ich erste Zwischenergebnisse haben. Und dementsprechend habe ich im Vorfeld schon mit Prompting ein Scoring-Modell bewertet auf Basis eines Sprachmodells oder entwickelt, ein Scoring-Modell entwickelt und habe einfach erstmals in meinem Leben versucht, diese neun Angebote in Form von PowerPoint-Folien durch dieses Scoring-Modell durchrauschen zu lassen und mir danach eine erste Einschätzung dieser neun Anbieter von der KI geben zu lassen.
  \end{quote}
  \textbf{Kategorie:} HK0 -- Übergreifende Kontextfaktoren\\
  \textbf{Weitere Kategorien:} HK1, HK2\\
  \textbf{Begr\"undung:} Die Passage beschreibt den organisationalen und zeitlichen Rahmen (Ausschreibung, Deadline, Anbieteranzahl), der den Kontext für den KI-Einsatz konstituiert und für das Verständnis der motivationalen Dynamik relevant ist.
  \\\textit{Memo: Die Textstelle ist ein klares Beispiel für proaktiven, kompetenzorientierten KI-Einsatz unter Zeitdruck. Der Ausdruck 'erstmals in meinem Leben' ist motivational bedeutsam und könnte im weiteren Verlauf auf ein Erleben von Neuheit oder Pioniercharakter hinweisen. Mehrfachcodierung HK1+HK2 ist hier gut begründbar; HK0 ergänzt den nötigen organisationalen Rahmen.}
\end{description}
\medskip

\begin{description}
  \item[\texttt{\#}\textbf{00:17:48}\texttt{\#} (Sinneinheit 20)]
  \begin{quote}
  Also für mich ist die KI immer so ein Thema, einerseits für die Recherche, also diese generative AI ist für mich so wichtig, aber natürlich auch diesen Automatisierungsanteil, den KI liefern kann. Also diese Mitarbeit im Team ist mir sehr, sehr wichtig. Und von dem her gesehen, früher hat es oft länger gedauert, kommt mir vor, bis man erste Ergebnisse für Besprechungen gehabt hat. Heutzutage kann ich eigentlich schon relativ schnell mit ersten Inputs einmal arbeiten. Und auch in der Weiterverarbeitung früher, nur ein Beispiel, wenn wir unsere Vorstände zum Beispiel zu den Bundesgremien entsendet haben, haben wir früher die Präsentationen im Haus an alle Fachbereiche verteilt und alle zum Beispiel acht Fachbereiche mussten zurückmelden, was in dieser Präsentation drinnen ist und welche Sichtweisen sie haben.
  \end{quote}
  \textbf{Kategorie:} HK0 -- Übergreifende Kontextfaktoren\\
  \textbf{Weitere Kategorien:} HK2\\
  \textbf{Begr\"undung:} Die Passage beschreibt den organisationalen Kontext der KI-Nutzung im Arbeitsalltag der Führungskraft, insbesondere den Einsatz generativer KI für Recherche und Automatisierung sowie bisherige Arbeitsprozesse im Team. Diese Schilderung dient als kontextueller Rahmen für das motivationale Erleben, ohne einem spezifischen Grundbedürfnis eindeutig zuordenbar zu sein.
  \\\textit{Memo: Die Sinneinheit bricht ab, bevor das konkrete Beispiel (Vorstände, Bundesgremien, Präsentationen) abgeschlossen wird. Der Vergleich 'früher–heute' ist ein klares Indiz für erlebten Kompetenz- und Effizienzgewinn. Die Textstelle enthält noch keinen expliziten Bezug zu Autonomie oder sozialer Eingebundenheit, obwohl der Verweis auf 'Mitarbeit im Team' als wichtig einen möglichen HK3-Anknüpfungspunkt für Folgepassagen andeutet.}
\end{description}
\medskip

\begin{description}
  \item[\texttt{\#}\textbf{00:17:48}\texttt{\#} (Sinneinheit 21)]
  \begin{quote}
  Mittlerweile machen wir zumindest diese erste Zusammenfassung, das erste Briefing mit Copilot und das auf SharePoint-Basis, damit es im gesicherten Bereich stattfindet. Aber du kannst den Vorstand viel besser vorbereiten, weil du dort acht bis zehn Präsentationen mit X-Folien gibst, einen Two-Pager mit, wo einmal eine Summary der Präsentation dabei ist und zusätzlich dann das Ergebnis der Fachbereichskollegen, das Feedback aus dem Fachbereich in Ergänzung. Und der Vorstand kann sich wirklich komprimiert einfach schnell vorbereiten. Das war früher mühsam. Und so, sage ich mal, ich könnte mir ehrlich gesagt die Arbeit ohne KI gar nicht mehr vorstellen. Wichtig ist mir halt, dass die Entscheidungsfindung trotzdem aber nur rein auf Basis valider Fakten und Einschätzungen von den Menschen getroffen werden.
  \end{quote}
  \textbf{Kategorie:} HK0 -- Übergreifende Kontextfaktoren\\
  \textbf{Weitere Kategorien:} HK1, HK2\\
  \textbf{Begr\"undung:} Der Hinweis auf die SharePoint-basierte, gesicherte Umgebung für die KI-Nutzung verweist auf organisationale und regulatorische Rahmenbedingungen im Bankensektor, die den Kontext des KI-Einsatzes strukturell einbetten.
  \\\textit{Memo: Die Sinneinheit enthält eine bemerkenswerte Spannungsdynamik: einerseits starke Abhängigkeit von und Begeisterung für KI ('könnte mir die Arbeit ohne KI gar nicht mehr vorstellen'), andererseits eine bewusste Grenzziehung hinsichtlich menschlicher Entscheidungshoheit. Diese Spannung zwischen KI-Nutzung und Erhalt humaner Kontrolle könnte für die Analyse des Autonomieerlebens im weiteren Verlauf besonders relevant sein.}
\end{description}
\medskip

\begin{description}
  \item[\texttt{\#}\textbf{00:23:20}\texttt{\#} (Sinneinheit 27)]
  \begin{quote}
  Für mich ist KI mittlerweile ein wesentlicher Bestandteil unseres täglichen, beziehungsweise zumindest wöchentlichen Tuns. Jeder muss sich in seinem Rollenprofil auch überlegen, wo er KI wirklich fix auch einbauen möchte, um seine Aufgaben zu erledigen. Jeder sollte bei mir im Team ganz klar überlegen, wie kann er KI nutzen, um effizienter, aber auch effektiver oder innovativer zu werden. Dementsprechend ist für mich auch mittlerweile, wenn dieses Thema Agentic AI bei uns wirklich hoffentlich stärker kommen sollte in den nächsten Jahren, muss man wirklich sich überlegen, wie man die Teams hybrid zusammensetzt.
  \end{quote}
  \textbf{Kategorie:} HK0 -- Übergreifende Kontextfaktoren\\
  \textbf{Weitere Kategorien:} HK1, HK3\\
  \textbf{Begr\"undung:} Die Passage beschreibt den organisationalen Rahmen der KI-Einführung als festen Bestandteil des Arbeitsalltags und skizziert den erwarteten zukünftigen Veränderungsdruck durch Agentic AI sowie die damit verbundene Notwendigkeit hybrider Teamzusammensetzungen. Dies stellt einen relevanten institutionellen Kontext für das motivationale Erleben dar, ohne einem spezifischen Grundbedürfnis eindeutig zugeordnet werden zu können.
  \\\textit{Memo: Die Sinneinheit verbindet Kontextbeschreibung (KI als etablierter Bestandteil) mit einer autonomieorientierten Führungshaltung und einem Ausblick auf veränderte Teamdynamiken. Besonders auffällig ist die normative Erwartungshaltung ('jeder muss', 'jeder sollte'), die eine gewisse Spannung zwischen proklamierter Autonomieförderung und implizitem Anpassungsdruck erzeugt – relevant für spätere Feincodierung.}
\end{description}
\medskip

\begin{description}
  \item[\texttt{\#}\textbf{00:24:20}\texttt{\#} (Sinneinheit 30)]
  \begin{quote}
  Und ich kann was verdienen damit. Vielleicht noch eine Rechnung. Wenn ich so einen Newsletter, das ist ein reiner Proof-of-Concept-Intern für steirische Firmen und Kundenberater, wo einfach die wesentlichsten steirischen, aber auch österreichischen Medien zu Wirtschaftsthemen abgegriffen werden und daraus branchenspezifische News einmal in der Woche, am Montag in der Früh ausgespielt werden, sowohl als Newsletter schriftlich, als aber auch als Podcast. Das kostet mir im ersten Schritt keine 14.000 Euro über externe Programmierung außerhalb unserer Systeme. Dann brauche ich noch das Hosting, das kostet 1.000 Euro, da bin ich bei 15.000 brutto.
  \end{quote}
  \textbf{Kategorie:} HK0 -- Übergreifende Kontextfaktoren\\
  \textbf{Weitere Kategorien:} HK2\\
  \textbf{Begr\"undung:} Die Passage beschreibt einen konkreten internen Proof-of-Concept mit KI-gestützter Newsletter- und Podcast-Produktion sowie die damit verbundenen Kostenkalkulationen. Sie liefert relevanten organisationalen und wirtschaftlichen Kontext für das Verständnis des KI-Einsatzes, ohne dass ein eindeutiger motivationaler Bezug zu einem der drei Grundbedürfnisse erkennbar ist.
  \\\textit{Memo: Die einleitende Aussage 'Ich kann was verdienen damit' könnte auf Autonomie- oder Kompetenzerleben hinweisen, bleibt aber zu vage für eine sichere HK1-Zuweisung. Der Hauptfokus liegt auf der technisch-wirtschaftlichen Projektbeschreibung. Möglicher Querverweis auf HK1, falls in benachbarten Sinneinheiten Eigeninitiative oder Entscheidungsfreiheit bei der Projektumsetzung thematisiert wird.}
\end{description}
\medskip

\begin{description}
  \item[\texttt{\#}\textbf{00:26:32}\texttt{\#} (Sinneinheit 34)]
  \begin{quote}
  Ich sage nur zum Beispiel das Thema Tokenisierung oder Coins, wenn das irgendwann einmal logischerweise mit monetären Kosten, mit einem monetären Aufwand verbunden ist, dann kann es uns auch passieren, dass das, was der Agent für uns macht, in Wahrheit teurer ist, weil es ein 0815-Thema ist, mit schlechter Qualität, als wenn es der Mensch weiterhin für uns gemacht hätte. Hört sich zwar ein bisschen perfide an, aber die nächste Herausforderung wird dann werden, ist dieser Agenteneinsatz wirklich effektiv und ist er effizient und wäre es nicht besser gewesen, das Thema wieder abzudrehen oder durch einen Menschen zu machen?
  \end{quote}
  \textbf{Kategorie:} HK0 -- Übergreifende Kontextfaktoren\\
  \textbf{Weitere Kategorien:} HK1, HK2\\
  \textbf{Begr\"undung:} Die Passage beschreibt organisationale Abwägungsprozesse rund um Effizienz und Kosteneffektivität beim KI-Agenteneinsatz und thematisiert damit strukturelle Rahmenbedingungen der KI-Einführung im institutionellen Kontext, ohne direkt auf ein Grundbedürfnis zu verweisen.
  \\\textit{Memo: Die Textstelle enthält eine nuancierte Kosten-Nutzen-Reflexion über KI-Agenten, die gleichzeitig Kontextfaktoren (Effizienz, Tokenisierungskosten), Autonomie (Entscheidung über Ab- oder Weiterbetrieb) und Kompetenz (Qualitätsvergleich Mensch vs. KI) berührt. Die Mehrfachkodierung erscheint gerechtfertigt, da alle drei Dimensionen substanziell angesprochen werden.}
\end{description}
\medskip

% ── HK1: Autonomieerleben ──────────────────────────────────────────────
\subsection*{HK1 -- Autonomieerleben}

\begin{description}
  \item[\texttt{\#}\textbf{00:02:18}\texttt{\#} (Sinneinheit 3)]
  \begin{quote}
  Es hat sich dann aber auch so entwickelt in weiterer Folge, dass das, was der KP-01 macht, sehr, sehr gut ist. Aber wir wollen jetzt so viele eigeninteressierte Mitarbeiterinnen und Mitarbeiter in unseren Bereichen haben, die sich auch dann selber Wissen aneignen. Wir haben dazu dann Schulungen auch aufgesetzt. M365, Co-Pilot, beiden Schulungen hat es auch bei uns im Haus gegeben. Einfach wirklich Fokus-Schulung zu Bau-API, aber immer mit KI-Kontext. Und wenn ich da jetzt sagen würde, Regeln gibt es logischerweise bitte auch, sowohl Richtung IT-Security und alles, was mit KI zu tun hat, ist bei uns geregelt im Haus. In der Praxis ist es dann so aber, dass viele Leute einfach Learning by Doing selber Sachen ausprobieren wollen. Und somit kann ich die Frage beantworten. Wir haben Regeln, wir haben Schulungen, wir haben es organisatorisch gelöst.
  \end{quote}
  \textbf{Kategorie:} HK1 -- Autonomieerleben\\
  \textbf{Weitere Kategorien:} HK0\\
  \textbf{Begr\"undung:} Die Interviewperson erwähnt explizit den Wunsch nach 'eigeninteressierten Mitarbeiterinnen und Mitarbeitern', die sich selbst Wissen aneignen, sowie das 'Learning by Doing' als bevorzugten Zugang der Mitarbeitenden. Dies verweist auf ein organisational angestrebtes und erlebtes Autonomieerleben beim Umgang mit KI-Tools, indem Selbststeuerung und intrinsische Exploration gefördert werden.
  \\\textit{Memo: Die Passage ist primär deskriptiv-kontextuell gehalten und schildert die organisationale KI-Einführungsarchitektur. Der Autonomiebezug ist indirekt, da er sich auf Mitarbeitende und nicht explizit auf die Führungskraft selbst bezieht. Für spätere Analysephasen relevant: Spannung zwischen institutionellen Regeln und dem Wunsch nach selbstgesteuertem Lernen (Autonomie vs. Kontrolle).}
\end{description}
\medskip

\begin{description}
  \item[\texttt{\#}\textbf{00:02:18}\texttt{\#} (Sinneinheit 4)]
  \begin{quote}
  Aber du kannst es dann, wenn es einmal Dynamik hat, glaube ich, nicht mehr so ganz aussteuern.
  \end{quote}
  \textbf{Kategorie:} HK1 -- Autonomieerleben\\
  \textbf{Begr\"undung:} Die Passage beschreibt den erlebten Kontrollverlust über einen einmal angestoßenen Prozess ('nicht mehr so ganz aussteuern'), was auf eine autonomiefrustrierende Erfahrung im Sinne eingeschränkter Handlungssteuerung hinweist.
  \\\textit{Memo: Die Textstelle ist ohne breiteren Kontext ambig – 'es' und 'Dynamik' sind nicht explizit benannt. Der Kern der Aussage verweist jedoch klar auf erlebten Steuerungsverlust, der am stärksten HK1 zuzuordnen ist. Kein expliziter KI-Bezug erkennbar, aber im Interviewkontext (motivationales Erleben von KI-Einsatz) motivational relevant.}
\end{description}
\medskip

\begin{description}
  \item[\texttt{\#}\textbf{00:04:52}\texttt{\#} (Sinneinheit 5)]
  \begin{quote}
  In meiner Funktion habe ich vorwiegend, also mein Hauptziel ist es immer, dass ich meine sechs Abteilungsleiter bestmöglich unterstütze und ihnen Rahmenbedingungen für Spitzenerfolge ermögliche, um es sehr hochtrabend zu sagen. Sprich, ich schaue, dass alle meine sechs Abteilungsleiter momentan nur Herren noch einfach bestmöglich ihre eigene Aufgabe und Verantwortung erfüllen können, dass sie dazu kompetent genug sind und dass, wenn es Probleme gibt, die ich für sie am besten regeln kann, dass sie dann sofort zu mir kommen, dass ich diese Probleme regle. Dadurch habe ich viele Besprechungen. Ich selber habe den Luxus, dass meine Leute so viel ohne mich machen, dass ich circa 50 Prozent meiner Zeit wirklich für persönliche Themen in Form von beruflichen Themen widmen kann, Mitarbeit in Projekten auf Bundes-, Landes-, aber auch Unternehmensebene hier in der Bank.
  \end{quote}
  \textbf{Kategorie:} HK1 -- Autonomieerleben\\
  \textbf{Weitere Kategorien:} HK0\\
  \textbf{Begr\"undung:} Die Interviewperson beschreibt explizit, dass sie rund 50 Prozent ihrer Zeit eigenständig für selbstgewählte berufliche Themen und Projekte nutzen kann – dies verweist auf einen hohen erlebten Handlungsspielraum und Selbstbestimmung in der Gestaltung der eigenen Arbeit.
  \\\textit{Memo: Die Sinneinheit enthält keinen direkten Bezug zu generativer KI, ist jedoch als Kontextpassage motivational relevant, da sie den Ausgangszustand der Autonomie und Führungsverantwortung beschreibt, vor dem spätere KI-bezogene Erfahrungen interpretiert werden müssen.}
\end{description}
\medskip

\begin{description}
  \item[\texttt{\#}\textbf{00:04:52}\texttt{\#} (Sinneinheit 6)]
  \begin{quote}
  Dementsprechend habe ich dann selber sehr viel mit der ganzen M365-Kollaborationssoftware zu tun. Und was ich immer früher gehabt habe in meinem Aufgabenbereich, brauche ich sehr oft Präsentationen. Präsentationen vor allem auch zu Themen, die es bei uns im Haus noch nicht gibt. Und ich bin immer an, ich versuche möglichst weit nach vorne zu denken. Und früher hätte ich da immer zwei, drei meiner Leute beauftragt, mir einmal das Thema zu recherchieren. Heutzutage kann ich da ganz offen und ehrlich sagen, mache ich das nur mit KI, mit klassisch, wenn es etwas ist, was ich außerhalb der Unternehmenssecurity machen darf, weil es ein ganz neues Thema ist, wie zum Beispiel Lebensphasenkonzepte im Vertrieb.
  \end{quote}
  \textbf{Kategorie:} HK1 -- Autonomieerleben\\
  \textbf{Weitere Kategorien:} HK2, HK3\\
  \textbf{Begr\"undung:} Die Interviewperson beschreibt, wie sie durch den Einsatz von KI eigenständig und ohne Delegation an Mitarbeitende Recherchen und Präsentationen erstellen kann. Dies verweist auf einen erweiterten Handlungsspielraum und größere Selbstbestimmtheit im Arbeitsprozess.
  \\\textit{Memo: Die Textstelle enthält einen relevanten Nebenbefund: Die explizite Erwähnung der 'Unternehmenssecurity' als Nutzungsbeschränkung könnte auch HK0 (regulatorisch-organisationaler Kontext) rechtfertigen, wurde jedoch hier nicht separat vergeben, da sie nur als Nebenbedingung erwähnt wird und der Fokus der Passage auf dem veränderten Arbeitsverhalten liegt. HK3 ist schwach ausgeprägt und könnte bei Revision als zu spekulativ eingestuft werden.}
\end{description}
\medskip

\begin{description}
  \item[\texttt{\#}\textbf{00:04:52}\texttt{\#} (Sinneinheit 7)]
  \begin{quote}
  Dann gehe ich ganz einmal beim Autofahren in Chat-GPT rein und mache das wirklich mit der Dialogfunktion, sprich mit der Konversationsfunktion und versuche das Thema mal so im Auto bei einer Wien-Fahrt zu interviewen. Und wenn das Ergebnis für mich noch stimmig ist, kann ich auch immer dann theoretisch am Desktop, am Notebook, wo ich immer das Thema vertiefe. Und dann lasse ich mal, das ist momentan ganz offen noch ein sehr schlampiger Weg, das Ergebnis nachher mit Notebook LM in eine Präsentation gießen, automatisiert. Die Präsentation ist in den meisten Fällen ja eh nur ein PDF mit Screenshots als Folien. Und das ist ehrlich gesagt oft ein brauchbares Ergebnis, aber nicht das, was ich mir vorstelle in weiterer Folge.
  \end{quote}
  \textbf{Kategorie:} HK1 -- Autonomieerleben\\
  \textbf{Weitere Kategorien:} HK2\\
  \textbf{Begr\"undung:} Die Passage beschreibt, wie die Interviewperson generative KI flexibel und selbstbestimmt in ihren persönlichen Arbeitsablauf integriert (Autofahrt, Desktop, Notebook) und dabei eigenständig über Zeitpunkt, Ort und Methode der Nutzung entscheidet. Gleichzeitig artikuliert sie eine kritische Distanz zum aktuellen Ergebnis ('schlampiger Weg', 'nicht das, was ich mir vorstelle'), was auf einen wahrgenommenen Handlungsspielraum und Gestaltungsanspruch hinweist.
  \\\textit{Memo: Die Sinneinheit zeigt eine pragmatisch-experimentelle Nutzungspraxis mit gleichzeitiger kritischer Selbstbewertung. Die Spannung zwischen 'brauchbarem Ergebnis' und dem eigenen Anspruchsniveau könnte auf eine latente Kompetenzfrustration hinweisen, auch wenn diese nicht explizit als Frustration formuliert wird. HK0 wurde nicht vergeben, da der motivationale Bezug hinreichend direkt ist.}
\end{description}
\medskip

\begin{description}
  \item[\texttt{\#}\textbf{00:06:27}\texttt{\#} (Sinneinheit 8)]
  \begin{quote}
  Wenn es Themen sind, wo ich selber den Sachverhalt noch nicht so im Detail verstehe und eher allgemeinere Themen, wo es erste Entscheidungen benötigt, dann nehme ich sehr oft KI zur Unterstützung dazu. Wenn es interne Themen sind, die in meiner Fachlichkeit gelagert sind, dann brauche ich dann in den meisten Fällen keine KI. Und eines möchte ich dir dazu sagen, das ist mir ganz wichtig, wenn ich KI-Unterstützung nehme, dann nur bei Themen, die ich fachlich wirklich nachvollziehen kann, weil das ist uns allen schon passiert, dass ein KI-Ergebnis, ein generiertes, einfach fachlich nicht schlüssig war, wenn man es im Detail geprüft hat. Ich nehme vor allem dann KI als Entscheidungshilfe, wenn es davor noch eine Recherche benötigt, wo ich Zusatzinformationen brauche, um mir meiner persönlichen Entscheidung sicherer zu werden.
  \end{quote}
  \textbf{Kategorie:} HK1 -- Autonomieerleben\\
  \textbf{Weitere Kategorien:} HK2\\
  \textbf{Begr\"undung:} Die Interviewperson beschreibt eine selbstbestimmte, selektive Nutzungsstrategie: KI wird bewusst nur in bestimmten Kontexten eingesetzt, während eigene fachliche Urteile Vorrang haben. Dies zeigt einen aktiv gestalteten Handlungsspielraum und die Wahrung von Entscheidungsfreiheit gegenüber KI-Outputs.
  \\\textit{Memo: Die Sinneinheit zeigt eine enge Verschränkung von Autonomie- und Kompetenzerleben: Die selbstbestimmte Nutzungsstrategie (HK1) ist funktional an das Kompetenzerleben (HK2) gekoppelt – Autonomie wird nur dort ausgeübt, wo eigene Fachlichkeit als Kontrollinstanz verfügbar ist. Bemerkenswert ist die explizite Warnung vor unkritischer KI-Nutzung als implizite Abgrenzung professioneller Identität.}
\end{description}
\medskip

\begin{description}
  \item[\texttt{\#}\textbf{00:07:54}\texttt{\#} (Sinneinheit 9)]
  \begin{quote}
  Ganz schräg, ich kann mittlerweile nicht mehr differenzieren, ob man gegenüber dieser KI, egal ob ich mit ihr spreche, mit ihr schreibe, ein Mensch ist oder kein Mensch ist. Was meine ich damit? Mir ist das ehrlich gesagt echt egal. Ich finde es, das ist ganz, für mich echt irritierend oft, ich bin ein sehr digitalaffiner Mensch und ich tue mir oft leichter, sogar mit der KI zu kommunizieren. Warum? Weil es ist oft so mühsam, einen Kollegen zum Zeitpunkt erst für einen Termin zu bekommen oder ihn jetzt anzurufen und er ist eh nicht erreichbar. Die KI ist für mich immer da, das ist für mich ein Convenience-Thema und ich kriege teilweise von einer KI Feedbacks oder Meinungen zugespiegelt, da muss ich vorher natürlich achten, wo kommt die Recherche her.
  \end{quote}
  \textbf{Kategorie:} HK1 -- Autonomieerleben\\
  \textbf{Weitere Kategorien:} HK3\\
  \textbf{Begr\"undung:} Die Passage beschreibt, wie die ständige Verfügbarkeit der KI den Handlungsspielraum der Interviewperson erweitert: Sie ist nicht mehr von Terminen oder Erreichbarkeit von Kollegen abhängig und kann selbstbestimmt und jederzeit auf Unterstützung zugreifen. Dies entspricht einer autonomieförderlichen Erfahrung durch erweiterte Gestaltungsfreiheit.
  \\\textit{Memo: Interessante Ambiguität: Die erlebte 'Erleichterung' im Umgang mit KI gegenüber Kollegen könnte sowohl als Autonomiegewinn (HK1) als auch als möglicher Rückzug aus sozialen Beziehungen (HK3, potenziell frustrierend) gelesen werden. Die Passage ist motivational bedeutsam, da sie auf eine subjektive Neubewertung sozialer Interaktion hinweist. Die Aussage 'mir ist es egal, ob Mensch oder nicht' deutet auf eine bemerkenswerte Verschiebung in der wahrgenommenen Qualität von Beziehungen hin.}
\end{description}
\medskip

\begin{description}
  \item[\texttt{\#}\textbf{00:07:54}\texttt{\#} (Sinneinheit 10)]
  \begin{quote}
  Ich versuche natürlich immer einen Prompt so darzustellen, dass der Prompt nur mehr valide Quellen berücksichtigt und so weiter mit Quellenangaben. Aber für mich ist die KI mittlerweile bei vielen Themen ein gleichwertiger Arbeitskollege oder Arbeitspartner und hat den großen Charme, dass ich manchen physischen Kollegen nicht mehr brauche. Ich liebe Menschen. Aber bei meinen Themen ist es oft lustiger, mit der KI zu arbeiten.
  \end{quote}
  \textbf{Kategorie:} HK1 -- Autonomieerleben\\
  \textbf{Weitere Kategorien:} HK2, HK3\\
  \textbf{Begr\"undung:} Die Interviewperson beschreibt, dass sie durch die KI-Nutzung selbstbestimmt entscheiden kann, mit wem sie zusammenarbeitet, und sich von bestimmten physischen Kollegen unabhängig macht. Dies verweist auf eine durch KI erweiterte Handlungsfreiheit und reduzierte interpersonelle Abhängigkeiten.
  \\\textit{Memo: Bemerkenswerte Ambivalenz: Die Interviewperson betont 'Ich liebe Menschen', relativiert dies aber unmittelbar durch die Präferenz für KI-Interaktion. Die KI wird explizit als sozialer Akteur ('Arbeitskollege') konzeptualisiert – dies könnte auf ein ungewöhnlich hohes Maß an parasozialer Bindung an die KI hindeuten, das im weiteren Verlauf des Interviews zu beobachten wäre.}
\end{description}
\medskip

\begin{description}
  \item[\texttt{\#}\textbf{00:10:14}\texttt{\#} (Sinneinheit 11)]
  \begin{quote}
  Wir hatten ja eine große Ausschreibung an externe Partner und hatten dann neun Angebote vorliegen. Ich bin von meinem Urlaub an dem Tag zurückgekommen, wo die Deadline war, für diese neun sehr hochwertigen Angebote. Es waren über 566 Folien, die ich da von diesen neun Anbietern bekommen habe. Und ich habe da im Vorfeld, weil ich ja gewusst habe, das ist erstens einmal zeitlich ein sehr, sehr knappes Thema. Freitag war die Deadline. Montag musste ich erste Zwischenergebnisse haben. Und dementsprechend habe ich im Vorfeld schon mit Prompting ein Scoring-Modell bewertet auf Basis eines Sprachmodells oder entwickelt, ein Scoring-Modell entwickelt und habe einfach erstmals in meinem Leben versucht, diese neun Angebote in Form von PowerPoint-Folien durch dieses Scoring-Modell durchrauschen zu lassen und mir danach eine erste Einschätzung dieser neun Anbieter von der KI geben zu lassen.
  \end{quote}
  \textbf{Kategorie:} HK1 -- Autonomieerleben\\
  \textbf{Weitere Kategorien:} HK0, HK2\\
  \textbf{Begr\"undung:} Die Interviewperson beschreibt eine eigeninitiative, selbstgesteuerte Entscheidung, ein Scoring-Modell via Prompting zu entwickeln und KI zur Angebotsbewertung einzusetzen — dies verweist auf erlebten Handlungsspielraum und selbstbestimmtes Vorgehen trotz extremem Zeitdruck.
  \\\textit{Memo: Die Textstelle ist ein klares Beispiel für proaktiven, kompetenzorientierten KI-Einsatz unter Zeitdruck. Der Ausdruck 'erstmals in meinem Leben' ist motivational bedeutsam und könnte im weiteren Verlauf auf ein Erleben von Neuheit oder Pioniercharakter hinweisen. Mehrfachcodierung HK1+HK2 ist hier gut begründbar; HK0 ergänzt den nötigen organisationalen Rahmen.}
\end{description}
\medskip

\begin{description}
  \item[\texttt{\#}\textbf{00:10:14}\texttt{\#} (Sinneinheit 12)]
  \begin{quote}
  Und ich muss offen sagen, niemand von uns kann 566 Folien von neun Anbietern so im Detail durchgehen in kurzer Zeit und verstehen, selbst wenn ich es aufgeteilt hätte auf höchst geschätzte Sektorkollegen. Und wie gesagt, jeder von uns macht drei dieser Angebote durch. Wir treffen uns dann in einer Teams und tauschen uns anhand eines Scoring-Modells, das von mir aus mit und ohne Karriere entwickelt worden ist, dazu strukturiert aus. Also ich muss echt sagen, das erste Ergebnis, was da rausgekommen ist, war für mich wirklich so gut brauchbar, dass man eine Shortlist, und das war auch das Ziel, aus diesen neuen Anbietern machen hätte können.
  \end{quote}
  \textbf{Kategorie:} HK1 -- Autonomieerleben\\
  \textbf{Weitere Kategorien:} HK2\\
  \textbf{Begr\"undung:} Die Interviewperson beschreibt, wie sie das KI-gestützte Scoring-Modell eigenverantwortlich mitentwickelt hat und damit den Entscheidungsprozess selbst strukturiert — dies verweist auf erlebte Gestaltungsfreiheit und Selbstbestimmung im Auswahlprozess.
  \\\textit{Memo: Die Textstelle illustriert ein klassisches KI-als-Kapazitätserweiterung-Erleben: Die Person erlebt sich nicht als ersetzt, sondern als in ihrer Kompetenz gestärkt und in ihrer Autonomie bestätigt, da sie das Bewertungsrahmenwerk selbst mitgestaltet hat. Mögliche Querverbindung zu HK3 (kollaborative Austauschrunde mit Sektorkollegen via Teams), die hier jedoch eher als organisationales Verfahren denn als Eingebundenheitserfahrung beschrieben wird.}
\end{description}
\medskip

\begin{description}
  \item[\texttt{\#}\textbf{00:12:58}\texttt{\#} (Sinneinheit 15)]
  \begin{quote}
  Und da habe ich wirklich gesehen, dass Spalten teilweise komplett, obwohl sie in beiden Listen drinnen waren, weggelassen worden sind, somit dann die Gesamtergebnisse nicht mehr gestimmt haben. Und wenn ich das nicht im Detail alles wirklich quasi fast mit einem Leuchtstift durchkontrolliert hätte gedanklich, wäre ich komplett über dieses Ergebnis und über diese Auswertungen und Empfehlungen gestolpert. Und das ist für mich der größte Knackpunkt. Ich traue mir ehrlich gesagt bei solchen Themen, wenn es um detaillierte Auswertungen, KI-gestützte geht, nicht über die KI drüber. Ich habe manch andere Themen schon probiert. So simpler Büroalltag bei mir ist es, dass ich alle drei Monate, einmal im Quartal, so ein Minigolf-Turnier für unser Haus organisiere. Das tue ich für den Betriebsrat. Und da kriege ich diese Anmeldungen alle per Automaten-Mail.
  \end{quote}
  \textbf{Kategorie:} HK1 -- Autonomieerleben\\
  \textbf{Weitere Kategorien:} HK2\\
  \textbf{Begr\"undung:} Die Interviewperson beschreibt, dass sie trotz KI-Unterstützung alle Ergebnisse eigenständig und detailliert kontrollieren muss, was auf einen eingeschränkten Vertrauensspielraum gegenüber der KI hinweist. Die Aussage 'Ich traue mir nicht über die KI drüber' zeigt, dass das selbstbestimmte Urteils- und Entscheidungsvermögen durch die Unzuverlässigkeit der KI eingeengt wird.
  \\\textit{Memo: Die Textstelle zeigt eine ambivalente Haltung: Einerseits wird die KI für bestimmte Aufgaben (komplexe Auswertungen) als unzuverlässig und fehleranfällig erlebt, andererseits wird ihr Einsatz für einfachere Routineaufgaben (Minigolf-Turnier, Büroalltag) angedeutet. Die Unterscheidung zwischen Aufgabentypen könnte in späteren Sinneinheiten als differenzierende Nutzungsstrategie weiter relevant werden.}
\end{description}
\medskip

\begin{description}
  \item[\texttt{\#}\textbf{00:12:58}\texttt{\#} (Sinneinheit 16)]
  \begin{quote}
  Wenn jemand ein Formular befüllt, kriege ich dazu das E-Mail mit der Anmeldung. Und da habe ich mir gedacht, Copilot-Routine kann ja nicht schwer sein, bastle mir aus all diesen E-Mails einfach eine Excel-Anmeldeliste. Und jedes Mal, wenn ich auf Knopfdruck, weil ein neues E-Mail da ist, schaut diese Anmeldeliste anders aus. Das ist jetzt aber auch kein Fehler der KI an dieser Stelle, möchte ich sagen, sondern ist mit Sicherheit ein Bedienfehler, weil ich das einfach zu schlecht geprompted habe. Zusammengefasst, alles, was in die detailliertere Analyse geht oder wo daraus wirklich konkrete Handlungen ausgelöst werden sollen, bin ich noch immer KI-skeptisch, auf Basis meiner mangelnden Kompetenz KI zu bedienen. Bei so Recherchethemen, wo es darum geht, neue Ideen zu entwickeln, Strategien einmal anzudenken, da liebe ich mittlerweile KI.
  \end{quote}
  \textbf{Kategorie:} HK1 -- Autonomieerleben\\
  \textbf{Weitere Kategorien:} HK2\\
  \textbf{Begr\"undung:} Die Passage beschreibt den selbstgesteuerten Versuch, eine eigene Copilot-Routine zu entwickeln, was auf einen wahrgenommenen Gestaltungsspielraum hinweist. Die Skepsis gegenüber KI bei konkreten Handlungsauslösungen deutet jedoch auf eine erlebte Einschränkung des autonomen Entscheidens hin, wenn die KI-Outputs nicht verlässlich kontrollierbar sind.
  \\\textit{Memo: Die Textstelle zeigt eine differenzierte Selbsteinschätzung: Die Person unterscheidet klar zwischen Anwendungsbereichen, in denen KI als kompetenzförderlich erlebt wird (Ideenentwicklung, Recherche), und solchen, in denen Kompetenzdefizite im Umgang mit KI zu Frustration führen. Die Selbstattribution des Fehlers ('Bedienfehler') ist motivational relevant, da sie Kontrollüberzeugungen und Kompetenzwahrnehmung berührt.}
\end{description}
\medskip

\begin{description}
  \item[\texttt{\#}\textbf{00:14:36}\texttt{\#} (Sinneinheit 17)]
  \begin{quote}
  Also dadurch, dass ich momentan noch so plump auftrete in Form dessen, dass das KI-Ergebnis, das mehrfach geprompte und iterierte, dann im Notebook LM einfach in so eine sehr bunte PowerPoint oder bunte PDF, muss ich sagen, gegossen wird, sage ich das immer gleich zu Beginn bei meinen Präsentationen dazu. Das Thema, wo ich immer sehr schnell versuche, mit KI erst zu recherchieren, das Ergebnis ist ja bitte über Notebook LM entstanden, die Visualisierung. Also ich würde nie verschweigen, dass ich KI verwendet habe dazu. Das ist mir sogar ganz, ganz wichtig, auch wenn es vielleicht beim Auditorium ein wenig die Skepsis erhöht. Aber in den meisten Fällen, wo ich mit KI jetzt aufgetreten bin und gerade die letzten zwei Tage in Salzburg auch mit in Summe zwei Präsentationen, hat sich niemand dadurch gestört gefühlt, dass das KI-aufbereitet war.
  \end{quote}
  \textbf{Kategorie:} HK1 -- Autonomieerleben\\
  \textbf{Weitere Kategorien:} HK3\\
  \textbf{Begr\"undung:} Die Passage beschreibt eine selbstbestimmte Entscheidung der Interviewperson, den KI-Einsatz transparent zu kommunizieren – auch gegen potenzielle Widerstände. Das aktive Gestalten der eigenen Präsentationspraxis und die bewusste Wahl zur Offenlegung verweisen auf erlebten Handlungsspielraum und Selbstbestimmung.
  \\\textit{Memo: Die Transparenz-Norm, die die Interviewperson für sich selbst setzt ('Das ist mir sogar ganz, ganz wichtig'), könnte auch auf Kompetenzerleben (professionelle Integrität) hinweisen. Die Passage bleibt jedoch primär im Bereich sozialer Akzeptanz und autonomer Werteorientierung; HK2 wird konservativ nicht vergeben.}
\end{description}
\medskip

\begin{description}
  \item[\texttt{\#}\textbf{00:14:36}\texttt{\#} (Sinneinheit 18)]
  \begin{quote}
  Dann, dass es mit Notebook LM eher unterschiedlich ausschaut und nicht mit unserem Master gemacht wurde, war auch nicht negativ. Und dadurch, sage ich mal, geht es mir gut damit. Wichtig ist mir, dass ich dadurch ja keine konkreten Entscheidungen dann losgeeist habe, sondern es war wirklich nur für strategische Überlegungen, sondern für einen Erst-Input gedacht.
  \end{quote}
  \textbf{Kategorie:} HK1 -- Autonomieerleben\\
  \textbf{Begr\"undung:} Die Interviewperson betont explizit, dass sie durch den KI-Einsatz keine konkreten Entscheidungen delegiert hat, sondern die KI nur als Erst-Input für strategische Überlegungen nutzt. Dies beschreibt eine bewusste Wahrung des eigenen Entscheidungsspielraums und damit autonomieerhaltende Handlungssteuerung.
  \\\textit{Memo: Die Formulierung 'keine konkreten Entscheidungen losgeeist' verweist auf eine klare Abgrenzung zwischen KI-unterstützter Reflexion und eigenverantwortlichem Entscheiden. Die positive emotionale Bilanz ('geht es mir gut damit') könnte auf Autonomiezufriedenheit hindeuten, ist aber für HK2 oder HK3 nicht ausreichend belegt.}
\end{description}
\medskip

\begin{description}
  \item[\texttt{\#}\textbf{00:17:48}\texttt{\#} (Sinneinheit 19)]
  \begin{quote}
  Ja, ich meine, wir haben ja bis vor ein, zwei Jahren ja alles ohne KI machen dürfen und müssen. Mir kommt vor, dass es früher teilweise mühsamer war. Und ich merke jetzt den Unterschied zwischen den Teammitgliedern, die mit KI arbeiten und den Mitgliedern, die weiterhin ohne KI arbeiten, aus welchen Gründen auch immer. Um mal schnell in die Gänge bei einem neuen Thema zu kommen, und bei mir sind oft eher neuere Themen am Tisch, ist KI mittlerweile für mich essentiell. Und ich gehe mittlerweile schon so weit, dass ich alle meine Kolleginnen und Kollegen dazu anhalte, bewusst mit KI zu arbeiten, damit sie sich selber einerseits Arbeitszeit ersparen, also effizienter werden, aber andererseits auch neue Sichtweisen gewinnen im Zusammenspiel mit der KI.
  \end{quote}
  \textbf{Kategorie:} HK1 -- Autonomieerleben\\
  \textbf{Weitere Kategorien:} HK2, HK3\\
  \textbf{Begr\"undung:} Die Interviewperson beschreibt, dass sie aktiv entscheidet, KI einzusetzen und Kolleg:innen dazu anhält – dies verweist auf einen selbstbestimmten, proaktiven Umgang mit KI-Tools, der den eigenen Handlungsspielraum erweitert. Die vergleichende Einschätzung ('früher mühsamer') unterstreicht das Erleben von gewonnener Handlungsfreiheit.
  \\\textit{Memo: Die Sinneinheit ist thematisch dicht und verbindet mehrere Dimensionen: eigene Effizienzgewinne (HK2), selbstbestimmte Nutzungsentscheidung und Multiplikatorrolle (HK1), sowie Teamdifferenzierung und Führungsimpuls (HK3). Keine HK0-Zuweisung notwendig, da alle relevanten Aspekte spezifischeren Kategorien zugeordnet werden können.}
\end{description}
\medskip

\begin{description}
  \item[\texttt{\#}\textbf{00:17:48}\texttt{\#} (Sinneinheit 21)]
  \begin{quote}
  Mittlerweile machen wir zumindest diese erste Zusammenfassung, das erste Briefing mit Copilot und das auf SharePoint-Basis, damit es im gesicherten Bereich stattfindet. Aber du kannst den Vorstand viel besser vorbereiten, weil du dort acht bis zehn Präsentationen mit X-Folien gibst, einen Two-Pager mit, wo einmal eine Summary der Präsentation dabei ist und zusätzlich dann das Ergebnis der Fachbereichskollegen, das Feedback aus dem Fachbereich in Ergänzung. Und der Vorstand kann sich wirklich komprimiert einfach schnell vorbereiten. Das war früher mühsam. Und so, sage ich mal, ich könnte mir ehrlich gesagt die Arbeit ohne KI gar nicht mehr vorstellen. Wichtig ist mir halt, dass die Entscheidungsfindung trotzdem aber nur rein auf Basis valider Fakten und Einschätzungen von den Menschen getroffen werden.
  \end{quote}
  \textbf{Kategorie:} HK1 -- Autonomieerleben\\
  \textbf{Weitere Kategorien:} HK2, HK0\\
  \textbf{Begr\"undung:} Die Passage beschreibt, wie KI die Interviewperson von mühsamen Routineaufgaben entlastet und dadurch Handlungsspielraum für qualitativ hochwertigere Vorstandsvorbereitung schafft. Gleichzeitig betont die Person explizit, dass die Entscheidungsfindung beim Menschen verbleiben soll, was auf ein aktives Bewahren von Selbstbestimmung im Entscheidungsprozess hinweist.
  \\\textit{Memo: Die Sinneinheit enthält eine bemerkenswerte Spannungsdynamik: einerseits starke Abhängigkeit von und Begeisterung für KI ('könnte mir die Arbeit ohne KI gar nicht mehr vorstellen'), andererseits eine bewusste Grenzziehung hinsichtlich menschlicher Entscheidungshoheit. Diese Spannung zwischen KI-Nutzung und Erhalt humaner Kontrolle könnte für die Analyse des Autonomieerlebens im weiteren Verlauf besonders relevant sein.}
\end{description}
\medskip

\begin{description}
  \item[\texttt{\#}\textbf{00:18:49}\texttt{\#} (Sinneinheit 23)]
  \begin{quote}
  Nein, ich sage mal, das müsste durch seine eigenen Gedanken und Erfahrungen vor allem, die die KI ja gar nicht in unserem Haus haben kann, bei der Beurteilung von Sachverhalten ergänzt werden. Also für mich, ich weiß nicht, ob es der richtige Begriff ist, müsste jedes Ergebnis ein Hybrid-Ergebnis sein. Mensch und Maschine. Also rein mit einem maschinellen Ergebnis zu mir zu kommen, da muss ich mich echt hinterfragen, für was ich den Mitarbeiter noch brauche, weil das hätte ich dann selber auch machen können. Ich möchte nicht haben, dass meine Leute mit der KI, mit dem KI-Ergebnis in den Diskurs gehen, für sich selber. Und vielleicht mit ein, zwei anderen Kollegen auch noch mal kurz abprüfen, evaluieren. Und das Ergebnis könnten sie mir noch in kürzerer Zeit in einer höheren Wertigkeit, glaube ich, dann überliefern.
  \end{quote}
  \textbf{Kategorie:} HK1 -- Autonomieerleben\\
  \textbf{Weitere Kategorien:} HK2, HK3\\
  \textbf{Begr\"undung:} Die Passage beschreibt die Erwartung der Führungskraft, dass Mitarbeiter eigene Gedanken und Erfahrungen einbringen sollen – das Hybrid-Ergebnis steht für einen Handlungsspielraum, in dem menschliche Entscheidungsfreiheit gegenüber rein maschinellen Ergebnissen bewahrt wird. Das Fehlen eigener Beiträge wird als Bedrohung der Sinnhaftigkeit der Mitarbeiterrolle erlebt.
  \\\textit{Memo: Zentrale Passage zum Hybrid-Konzept: Die Führungskraft formuliert ein normatives Leitbild für den KI-Einsatz, das alle drei Grundbedürfnisse berührt. Besonders auffällig ist die implizite Sorge um Rollenentleerung und Expertiseentwertung (HK2), die eng mit der Frage nach dem Sinn menschlicher Mitarbeit verknüpft ist.}
\end{description}
\medskip

\begin{description}
  \item[\texttt{\#}\textbf{00:23:20}\texttt{\#} (Sinneinheit 27)]
  \begin{quote}
  Für mich ist KI mittlerweile ein wesentlicher Bestandteil unseres täglichen, beziehungsweise zumindest wöchentlichen Tuns. Jeder muss sich in seinem Rollenprofil auch überlegen, wo er KI wirklich fix auch einbauen möchte, um seine Aufgaben zu erledigen. Jeder sollte bei mir im Team ganz klar überlegen, wie kann er KI nutzen, um effizienter, aber auch effektiver oder innovativer zu werden. Dementsprechend ist für mich auch mittlerweile, wenn dieses Thema Agentic AI bei uns wirklich hoffentlich stärker kommen sollte in den nächsten Jahren, muss man wirklich sich überlegen, wie man die Teams hybrid zusammensetzt.
  \end{quote}
  \textbf{Kategorie:} HK1 -- Autonomieerleben\\
  \textbf{Weitere Kategorien:} HK0, HK3\\
  \textbf{Begr\"undung:} Die Interviewperson beschreibt, dass jedes Teammitglied eigenständig entscheiden soll, wie es KI in sein Rollenprofil integriert, um effizienter, effektiver oder innovativer zu werden. Dies verweist auf eine autonomieförderliche Gestaltung: Die Führungskraft eröffnet Handlungs- und Entscheidungsspielräume bei der KI-Nutzung, anstatt Vorgaben zu machen.
  \\\textit{Memo: Die Sinneinheit verbindet Kontextbeschreibung (KI als etablierter Bestandteil) mit einer autonomieorientierten Führungshaltung und einem Ausblick auf veränderte Teamdynamiken. Besonders auffällig ist die normative Erwartungshaltung ('jeder muss', 'jeder sollte'), die eine gewisse Spannung zwischen proklamierter Autonomieförderung und implizitem Anpassungsdruck erzeugt – relevant für spätere Feincodierung.}
\end{description}
\medskip

\begin{description}
  \item[\texttt{\#}\textbf{00:23:20}\texttt{\#} (Sinneinheit 28)]
  \begin{quote}
  Weil wenn ich dann wirklich Kolleginnen und Kollegen mit teilweise 20-jährigem Know-how, sprich Wissen und Erfahrung habe, plus das ganze Agenze durch schlaue Agenten, die einfach viele so bürokratische oder vielleicht auch Recherchetätigkeiten im Unternehmenswissen, von der Rechtsabteilung bis zum Datenschutz bis zum Marketing für uns erledigen können, dann kann er viel schlagkräftiger werden in unserem Haus. Und es gibt zum Beispiel für mich Beispiele, wo ich mit menschlichem Aufwand viel zu viele Kosten in die Hand nehmen würde, wenn ich das einmal schlau über AI-Netzwerke, über kleine mit Automatisierung löse, kann ich auf einmal zum Beispiel einen Newsletter für steirische Firmen und Kundenberater zu den wesentlichsten Wirtschafts-News der Steiermark kostengünstigst erwirken, ganzjährig.
  \end{quote}
  \textbf{Kategorie:} HK1 -- Autonomieerleben\\
  \textbf{Weitere Kategorien:} HK2\\
  \textbf{Begr\"undung:} Die Interviewperson beschreibt, wie KI-gestützte Automatisierung neue Handlungsspielräume eröffnet, die ohne KI aufgrund zu hoher Kosten nicht realisierbar wären – etwa der ganzjährige Newsletter für steirische Firmen. Dies entspricht einer autonomieförderlichen Erfahrung durch erweiterte Gestaltungsmöglichkeiten.
  \\\textit{Memo: Die Textstelle verbindet instrumentelle Effizienzgewinne mit einem klaren Erleben erweiterter Handlungs- und Wirksamkeitspotenziale. Die Metapher 'schlagkräftiger werden' ist besonders aufschlussreich für das Kompetenzerleben. Eine Zuordnung zu HK3 wäre denkbar (Teamsynergie Mensch-KI), erscheint jedoch nicht hinreichend explizit belegt.}
\end{description}
\medskip

\begin{description}
  \item[\texttt{\#}\textbf{00:23:20}\texttt{\#} (Sinneinheit 29)]
  \begin{quote}
  Früher, wenn ich das gemacht hätte über Manpower oder Womenpower in meinem Team, da wäre es mir zu schade gewesen, diese wertvollen Ressourcen für so einen Newsletter einzusetzen.
  \end{quote}
  \textbf{Kategorie:} HK1 -- Autonomieerleben\\
  \textbf{Weitere Kategorien:} HK2\\
  \textbf{Begr\"undung:} Die Passage beschreibt, wie die Interviewperson früher bei der Ressourcenallokation im Team abwägen musste und bestimmte Aufgaben bewusst nicht delegierte – dies verweist auf den erlebten Handlungsspielraum und die Entscheidungsfreiheit über den Einsatz von Teamkapazitäten.
  \\\textit{Memo: Die Sinneinheit gewinnt ihren eigentlichen motivationalen Gehalt erst im Kontrast zur KI-gestützten Gegenwart: Implizit wird eine frühere Einschränkung (kein Newsletter wegen Ressourcenknappheit) beschrieben, die durch KI-Einsatz aufgehoben wurde. Für die Vollständigkeit der Interpretation sollte der unmittelbare Kontext (Folgesätze) herangezogen werden.}
\end{description}
\medskip

\begin{description}
  \item[\texttt{\#}\textbf{00:24:20}\texttt{\#} (Sinneinheit 31)]
  \begin{quote}
  Und wenn ich es schaffe, dass jede steirische Raiffeisenbank in weiterer Folge von dieser Dienstleistung profitieren möchte und 50 Euro im Monat, egal wie viele Users sie hat, bezahlt, bin ich bei 40.000, 50.000 und 2.000 Euro, bin ich bei 24.000 Euro Erträgen, obwohl ich einmal im ersten Schritt circa 15.000 investiert habe und das im ersten Jahr. Und das sind Themen, die würde ich mit menschlichem Arbeitseinsatz nie profitabel gestalten können.
  \end{quote}
  \textbf{Kategorie:} HK1 -- Autonomieerleben\\
  \textbf{Weitere Kategorien:} HK2\\
  \textbf{Begr\"undung:} Die Interviewperson schildert, wie sie eigeninitiativ ein skalierbares Geschäftsmodell gestaltet und dabei Gestaltungsspielräume nutzt, die KI erst ermöglicht. Der selbstbestimmte unternehmerische Handlungsrahmen – von der Investitionsentscheidung bis zum Ertragsmodell – verweist auf autonomieförderliches Erleben.
  \\\textit{Memo: Die Textstelle enthält eine konkrete Kalkulation als Beleg für den wahrgenommenen KI-Mehrwert. Die Verbindung von Kompetenz- und Autonomieerleben ist hier eng verknüpft: Die KI erweitert nicht nur die Fähigkeiten, sondern eröffnet auch einen selbstgestalteten unternehmerischen Spielraum. Eine Zuordnung zu HK0 (institutioneller Kontext Raiffeisenbanken) wäre denkbar, tritt jedoch hinter die motivationale Dimension zurück.}
\end{description}
\medskip

\begin{description}
  \item[\texttt{\#}\textbf{00:26:32}\texttt{\#} (Sinneinheit 33)]
  \begin{quote}
  Ja, also mir ist wichtig, dass wir vor allem bestmöglich auch diese KI-Möglichkeiten gezielt einsetzen, eben von der Generierung bis zur Automatisierung und so weiter, aber trotzdem permanent Messpunkte entwickeln, um zu schauen, ob wir dadurch wirklich effizienter und effektiver werden. Und vor allem, wenn wirklich losgelöste, autonome Agenten dann für uns Tätigkeiten machen, wie können wir wirklich permanent bewerten, ob das, was diese Zusatzkollegen für uns generieren und produzieren, wirklich in das, was wir erreichen wollen, die strategischen Ziele, wirklich auch einzahlt. Wenn es nur operative Entlastung ist, ist es auch gut, aber es muss eine gewisse Gratwanderung sein, bis wohin hilft es uns wirklich weiter und wo ist vielleicht zu viel des Guten.
  \end{quote}
  \textbf{Kategorie:} HK1 -- Autonomieerleben\\
  \textbf{Weitere Kategorien:} HK2\\
  \textbf{Begr\"undung:} Die Interviewperson beschreibt den Wunsch, Kontroll- und Bewertungsinstanzen ('Messpunkte') gegenüber autonomen KI-Agenten zu etablieren, um den eigenen Handlungsspielraum und die strategische Steuerungsfähigkeit zu erhalten. Das Bild der 'Gratwanderung' verweist auf das Spannungsfeld zwischen KI-Delegation und selbstbestimmter Kontrolle über Prozesse und Ziele.
  \\\textit{Memo: Die Sinneinheit enthält eine charakteristische Ambivalenz: KI-Nutzung wird begrüßt (operative Entlastung), gleichzeitig wird die Notwendigkeit eigenständiger Steuerungs- und Bewertungskompetenz betont. Die Formulierung 'Zusatzkollegen' ist semantisch auffällig und könnte in Bezug auf HK3 (soziale Eingebundenheit, Mensch-KI-Beziehung) in späteren Durchgängen weiter reflektiert werden.}
\end{description}
\medskip

\begin{description}
  \item[\texttt{\#}\textbf{00:26:32}\texttt{\#} (Sinneinheit 34)]
  \begin{quote}
  Ich sage nur zum Beispiel das Thema Tokenisierung oder Coins, wenn das irgendwann einmal logischerweise mit monetären Kosten, mit einem monetären Aufwand verbunden ist, dann kann es uns auch passieren, dass das, was der Agent für uns macht, in Wahrheit teurer ist, weil es ein 0815-Thema ist, mit schlechter Qualität, als wenn es der Mensch weiterhin für uns gemacht hätte. Hört sich zwar ein bisschen perfide an, aber die nächste Herausforderung wird dann werden, ist dieser Agenteneinsatz wirklich effektiv und ist er effizient und wäre es nicht besser gewesen, das Thema wieder abzudrehen oder durch einen Menschen zu machen?
  \end{quote}
  \textbf{Kategorie:} HK1 -- Autonomieerleben\\
  \textbf{Weitere Kategorien:} HK0, HK2\\
  \textbf{Begr\"undung:} Die Interviewperson beschreibt die Möglichkeit, den Agenteneinsatz aktiv zu evaluieren und gegebenenfalls 'abzudrehen' oder durch menschliche Arbeit zu ersetzen — dies verweist auf einen wahrgenommenen Entscheidungsspielraum und damit auf das Erleben von Handlungsautonomie gegenüber der KI-Nutzung.
  \\\textit{Memo: Die Textstelle enthält eine nuancierte Kosten-Nutzen-Reflexion über KI-Agenten, die gleichzeitig Kontextfaktoren (Effizienz, Tokenisierungskosten), Autonomie (Entscheidung über Ab- oder Weiterbetrieb) und Kompetenz (Qualitätsvergleich Mensch vs. KI) berührt. Die Mehrfachkodierung erscheint gerechtfertigt, da alle drei Dimensionen substanziell angesprochen werden.}
\end{description}
\medskip

\begin{description}
  \item[\texttt{\#}\textbf{00:26:54}\texttt{\#} (Sinneinheit 35)]
  \begin{quote}
  Konsumation ist cool, aber sie muss trotzdem gezielt strukturiert erfolgen und man darf sich einfach nicht zu sehr auf alle neuen Möglichkeiten verlassen und dann sich wundern, wenn alles aus dem Ruder läuft.
  \end{quote}
  \textbf{Kategorie:} HK1 -- Autonomieerleben\\
  \textbf{Begr\"undung:} Die Passage beschreibt die Notwendigkeit, den eigenen Umgang mit KI-Tools bewusst und selbstgesteuert zu gestalten, um Kontrollverlust zu vermeiden. Dies verweist auf das Erleben von Handlungsspielraum und Selbstbestimmung im Umgang mit generativer KI — konkret die eigenverantwortliche Steuerung statt unreflektierter Abhängigkeit.
  \\\textit{Memo: Die Formulierung 'sich nicht zu sehr verlassen' deutet auf eine wahrgenommene Gefährdung des Autonomieerlebens durch unkritische KI-Nutzung hin. Es schwingt eine präventive Haltung gegenüber autonomiefrustrierenden Dynamiken (Kontrollverlust, Abhängigkeit) mit, ohne dass diese bereits eingetreten wären.}
\end{description}
\medskip

% ── HK2: Kompetenzerleben ──────────────────────────────────────────────
\subsection*{HK2 -- Kompetenzerleben}

\begin{description}
  \item[\texttt{\#}\textbf{00:04:52}\texttt{\#} (Sinneinheit 6)]
  \begin{quote}
  Dementsprechend habe ich dann selber sehr viel mit der ganzen M365-Kollaborationssoftware zu tun. Und was ich immer früher gehabt habe in meinem Aufgabenbereich, brauche ich sehr oft Präsentationen. Präsentationen vor allem auch zu Themen, die es bei uns im Haus noch nicht gibt. Und ich bin immer an, ich versuche möglichst weit nach vorne zu denken. Und früher hätte ich da immer zwei, drei meiner Leute beauftragt, mir einmal das Thema zu recherchieren. Heutzutage kann ich da ganz offen und ehrlich sagen, mache ich das nur mit KI, mit klassisch, wenn es etwas ist, was ich außerhalb der Unternehmenssecurity machen darf, weil es ein ganz neues Thema ist, wie zum Beispiel Lebensphasenkonzepte im Vertrieb.
  \end{quote}
  \textbf{Kategorie:} HK2 -- Kompetenzerleben\\
  \textbf{Weitere Kategorien:} HK1, HK3\\
  \textbf{Begr\"undung:} Die Passage zeigt, dass die Interviewperson durch KI neue Themengebiete eigenständig erschließen kann, was auf eine Erweiterung der eigenen Wirksamkeit und Handlungsfähigkeit hinweist – früher war dafür die Beauftragung mehrerer Mitarbeitender notwendig.
  \\\textit{Memo: Die Textstelle enthält einen relevanten Nebenbefund: Die explizite Erwähnung der 'Unternehmenssecurity' als Nutzungsbeschränkung könnte auch HK0 (regulatorisch-organisationaler Kontext) rechtfertigen, wurde jedoch hier nicht separat vergeben, da sie nur als Nebenbedingung erwähnt wird und der Fokus der Passage auf dem veränderten Arbeitsverhalten liegt. HK3 ist schwach ausgeprägt und könnte bei Revision als zu spekulativ eingestuft werden.}
\end{description}
\medskip

\begin{description}
  \item[\texttt{\#}\textbf{00:04:52}\texttt{\#} (Sinneinheit 7)]
  \begin{quote}
  Dann gehe ich ganz einmal beim Autofahren in Chat-GPT rein und mache das wirklich mit der Dialogfunktion, sprich mit der Konversationsfunktion und versuche das Thema mal so im Auto bei einer Wien-Fahrt zu interviewen. Und wenn das Ergebnis für mich noch stimmig ist, kann ich auch immer dann theoretisch am Desktop, am Notebook, wo ich immer das Thema vertiefe. Und dann lasse ich mal, das ist momentan ganz offen noch ein sehr schlampiger Weg, das Ergebnis nachher mit Notebook LM in eine Präsentation gießen, automatisiert. Die Präsentation ist in den meisten Fällen ja eh nur ein PDF mit Screenshots als Folien. Und das ist ehrlich gesagt oft ein brauchbares Ergebnis, aber nicht das, was ich mir vorstelle in weiterer Folge.
  \end{quote}
  \textbf{Kategorie:} HK2 -- Kompetenzerleben\\
  \textbf{Weitere Kategorien:} HK1\\
  \textbf{Begr\"undung:} Die Interviewperson reflektiert explizit die Qualität des KI-generierten Outputs ('oft ein brauchbares Ergebnis, aber nicht das, was ich mir vorstelle') und setzt damit einen professionellen Qualitätsmaßstab an, der auf ein Spannungsverhältnis zwischen erweiterter Leistungsfähigkeit durch KI und dem eigenen Kompetenzanspruch hindeutet.
  \\\textit{Memo: Die Sinneinheit zeigt eine pragmatisch-experimentelle Nutzungspraxis mit gleichzeitiger kritischer Selbstbewertung. Die Spannung zwischen 'brauchbarem Ergebnis' und dem eigenen Anspruchsniveau könnte auf eine latente Kompetenzfrustration hinweisen, auch wenn diese nicht explizit als Frustration formuliert wird. HK0 wurde nicht vergeben, da der motivationale Bezug hinreichend direkt ist.}
\end{description}
\medskip

\begin{description}
  \item[\texttt{\#}\textbf{00:06:27}\texttt{\#} (Sinneinheit 8)]
  \begin{quote}
  Wenn es Themen sind, wo ich selber den Sachverhalt noch nicht so im Detail verstehe und eher allgemeinere Themen, wo es erste Entscheidungen benötigt, dann nehme ich sehr oft KI zur Unterstützung dazu. Wenn es interne Themen sind, die in meiner Fachlichkeit gelagert sind, dann brauche ich dann in den meisten Fällen keine KI. Und eines möchte ich dir dazu sagen, das ist mir ganz wichtig, wenn ich KI-Unterstützung nehme, dann nur bei Themen, die ich fachlich wirklich nachvollziehen kann, weil das ist uns allen schon passiert, dass ein KI-Ergebnis, ein generiertes, einfach fachlich nicht schlüssig war, wenn man es im Detail geprüft hat. Ich nehme vor allem dann KI als Entscheidungshilfe, wenn es davor noch eine Recherche benötigt, wo ich Zusatzinformationen brauche, um mir meiner persönlichen Entscheidung sicherer zu werden.
  \end{quote}
  \textbf{Kategorie:} HK2 -- Kompetenzerleben\\
  \textbf{Weitere Kategorien:} HK1\\
  \textbf{Begr\"undung:} Die Passage thematisiert explizit das Erfordernis fachlicher Nachvollziehbarkeit als Voraussetzung für KI-Nutzung und verweist auf Erfahrungen mit fachlich nicht schlüssigen KI-Ergebnissen. Dies reflektiert ein ausgeprägtes Kompetenzselbstbild, das als Korrektiv gegenüber KI-Outputs fungiert und die eigene Expertise als unverzichtbar markiert.
  \\\textit{Memo: Die Sinneinheit zeigt eine enge Verschränkung von Autonomie- und Kompetenzerleben: Die selbstbestimmte Nutzungsstrategie (HK1) ist funktional an das Kompetenzerleben (HK2) gekoppelt – Autonomie wird nur dort ausgeübt, wo eigene Fachlichkeit als Kontrollinstanz verfügbar ist. Bemerkenswert ist die explizite Warnung vor unkritischer KI-Nutzung als implizite Abgrenzung professioneller Identität.}
\end{description}
\medskip

\begin{description}
  \item[\texttt{\#}\textbf{00:07:54}\texttt{\#} (Sinneinheit 10)]
  \begin{quote}
  Ich versuche natürlich immer einen Prompt so darzustellen, dass der Prompt nur mehr valide Quellen berücksichtigt und so weiter mit Quellenangaben. Aber für mich ist die KI mittlerweile bei vielen Themen ein gleichwertiger Arbeitskollege oder Arbeitspartner und hat den großen Charme, dass ich manchen physischen Kollegen nicht mehr brauche. Ich liebe Menschen. Aber bei meinen Themen ist es oft lustiger, mit der KI zu arbeiten.
  \end{quote}
  \textbf{Kategorie:} HK2 -- Kompetenzerleben\\
  \textbf{Weitere Kategorien:} HK1, HK3\\
  \textbf{Begr\"undung:} Die bewusste Formulierung valider Prompts mit Quellenangaben verweist auf ein aktives Kompetenzmanagement im Umgang mit KI. Die Einschätzung der KI als 'gleichwertiger Arbeitspartner' deutet darauf hin, dass die Interviewperson sich selbst als kompetent genug erlebt, um mit der KI auf Augenhöhe zu arbeiten und qualitativ hochwertige Ergebnisse zu erzielen.
  \\\textit{Memo: Bemerkenswerte Ambivalenz: Die Interviewperson betont 'Ich liebe Menschen', relativiert dies aber unmittelbar durch die Präferenz für KI-Interaktion. Die KI wird explizit als sozialer Akteur ('Arbeitskollege') konzeptualisiert – dies könnte auf ein ungewöhnlich hohes Maß an parasozialer Bindung an die KI hindeuten, das im weiteren Verlauf des Interviews zu beobachten wäre.}
\end{description}
\medskip

\begin{description}
  \item[\texttt{\#}\textbf{00:10:14}\texttt{\#} (Sinneinheit 11)]
  \begin{quote}
  Wir hatten ja eine große Ausschreibung an externe Partner und hatten dann neun Angebote vorliegen. Ich bin von meinem Urlaub an dem Tag zurückgekommen, wo die Deadline war, für diese neun sehr hochwertigen Angebote. Es waren über 566 Folien, die ich da von diesen neun Anbietern bekommen habe. Und ich habe da im Vorfeld, weil ich ja gewusst habe, das ist erstens einmal zeitlich ein sehr, sehr knappes Thema. Freitag war die Deadline. Montag musste ich erste Zwischenergebnisse haben. Und dementsprechend habe ich im Vorfeld schon mit Prompting ein Scoring-Modell bewertet auf Basis eines Sprachmodells oder entwickelt, ein Scoring-Modell entwickelt und habe einfach erstmals in meinem Leben versucht, diese neun Angebote in Form von PowerPoint-Folien durch dieses Scoring-Modell durchrauschen zu lassen und mir danach eine erste Einschätzung dieser neun Anbieter von der KI geben zu lassen.
  \end{quote}
  \textbf{Kategorie:} HK2 -- Kompetenzerleben\\
  \textbf{Weitere Kategorien:} HK0, HK1\\
  \textbf{Begr\"undung:} Die Passage schildert, wie die Interviewperson durch den gezielten Einsatz von Prompting und einem selbst entwickelten Scoring-Modell eine komplexe fachliche Aufgabe bewältigt, was auf ein Erleben erweiterter Wirksamkeit und neuer Kompetenzdomänen im Umgang mit generativer KI hindeutet.
  \\\textit{Memo: Die Textstelle ist ein klares Beispiel für proaktiven, kompetenzorientierten KI-Einsatz unter Zeitdruck. Der Ausdruck 'erstmals in meinem Leben' ist motivational bedeutsam und könnte im weiteren Verlauf auf ein Erleben von Neuheit oder Pioniercharakter hinweisen. Mehrfachcodierung HK1+HK2 ist hier gut begründbar; HK0 ergänzt den nötigen organisationalen Rahmen.}
\end{description}
\medskip

\begin{description}
  \item[\texttt{\#}\textbf{00:10:14}\texttt{\#} (Sinneinheit 12)]
  \begin{quote}
  Und ich muss offen sagen, niemand von uns kann 566 Folien von neun Anbietern so im Detail durchgehen in kurzer Zeit und verstehen, selbst wenn ich es aufgeteilt hätte auf höchst geschätzte Sektorkollegen. Und wie gesagt, jeder von uns macht drei dieser Angebote durch. Wir treffen uns dann in einer Teams und tauschen uns anhand eines Scoring-Modells, das von mir aus mit und ohne Karriere entwickelt worden ist, dazu strukturiert aus. Also ich muss echt sagen, das erste Ergebnis, was da rausgekommen ist, war für mich wirklich so gut brauchbar, dass man eine Shortlist, und das war auch das Ziel, aus diesen neuen Anbietern machen hätte können.
  \end{quote}
  \textbf{Kategorie:} HK2 -- Kompetenzerleben\\
  \textbf{Weitere Kategorien:} HK1\\
  \textbf{Begr\"undung:} Die Passage schildert, dass die KI die Bewältigung einer kognitiv kaum handhabbaren Aufgabe (566 Folien von neun Anbietern) ermöglichte und ein qualitativ hochwertiges, direkt nutzbares Ergebnis (Shortlist) lieferte — dies verweist auf erlebte Erweiterung der eigenen Handlungsfähigkeit und Wirksamkeit.
  \\\textit{Memo: Die Textstelle illustriert ein klassisches KI-als-Kapazitätserweiterung-Erleben: Die Person erlebt sich nicht als ersetzt, sondern als in ihrer Kompetenz gestärkt und in ihrer Autonomie bestätigt, da sie das Bewertungsrahmenwerk selbst mitgestaltet hat. Mögliche Querverbindung zu HK3 (kollaborative Austauschrunde mit Sektorkollegen via Teams), die hier jedoch eher als organisationales Verfahren denn als Eingebundenheitserfahrung beschrieben wird.}
\end{description}
\medskip

\begin{description}
  \item[\texttt{\#}\textbf{00:10:14}\texttt{\#} (Sinneinheit 13)]
  \begin{quote}
  Wenn ich das gemacht hätte mit menschlichen Ressourcen, mit den unterschiedlichsten subjektiven Zugängen, die jeder in sich trägt, hätte es wahrscheinlich zwei Wochen gedauert und ich bin mir ehrlich gesagt nicht sicher, ob das Ergebnis, um eine Shortlist zu definieren, ein besseres dadurch geworden wäre.
  \end{quote}
  \textbf{Kategorie:} HK2 -- Kompetenzerleben\\
  \textbf{Begr\"undung:} Die Passage beschreibt, wie generative KI eine Aufgabe effizienter und in vergleichbarer oder besserer Qualität erledigt als ein Team menschlicher Ressourcen es könnte. Dies verweist auf eine erweiterte Wirksamkeit durch KI-Einsatz, was dem Kompetenzerleben zuzuordnen ist.
  \\\textit{Memo: Die Formulierung 'nicht sicher, ob das Ergebnis ein besseres dadurch geworden wäre' deutet auf eine ambivalente Bewertung hin: Die KI wird nicht als eindeutig überlegen dargestellt, sondern als mindestens gleichwertig bei deutlich geringerem Zeitaufwand. Dies könnte auch einen leichten Unterton von Expertise-Relativierung menschlicher Kollegen enthalten, bleibt aber primär im Bereich Kompetenzerleben (Effizienz und Ergebnisqualität).}
\end{description}
\medskip

\begin{description}
  \item[\texttt{\#}\textbf{00:12:58}\texttt{\#} (Sinneinheit 14)]
  \begin{quote}
  Ja, das ist spannend bei anderen Themen wiederum. Da bin ich aber wirklich selbst schuld, weil einfach mein KI-Level von der Kompetenz noch zu gering ist. Ich müsste wirklich da vertiefte Schulungen machen. Ich habe relativ oft noch immer Excel-Listen am Tisch, wo ich Ergebnisse von reifen Banken bewerten muss, um daraus einfach auf Basis dieser Analyse neue Ideen, Lösungen zu entwickeln. Und da habe ich einfach viel zu simpel gedacht, habe einfach diese Ergebnislisten im Excel hergenommen, der 40 steirischen Banken, habe es im Excel mit Copilot mir auswerten lassen, was sind die größten Auffälligkeiten, wenn du diese beiden Excel-Ergebnisse vergleichst und so weiter, und basteln wir daraus eine neue Excel, die das und das berücksichtigt, inklusive grafischer Darstellung.
  \end{quote}
  \textbf{Kategorie:} HK2 -- Kompetenzerleben\\
  \textbf{Begr\"undung:} Die Interviewperson bewertet ihr eigenes KI-Kompetenzniveau als unzureichend ('viel zu simpel gedacht') und benennt explizit Schulungsbedarf. Gleichzeitig beschreibt sie eine konkrete Anwendungssituation mit Copilot, in der sie KI-gestützt Analyseaufgaben bewältigt — was sowohl auf Kompetenzfrustrierung (zu geringe KI-Kompetenz) als auch auf kompetenzerweiternde Nutzung (Auswertung komplexer Daten, grafische Darstellung) verweist.
  \\\textit{Memo: Die Passage enthält eine bemerkenswerte Selbstattribution: Die Interviewperson macht sich selbst verantwortlich ('bin ich selbst schuld') für die suboptimale KI-Nutzung, was auf eine internale Kontrollüberzeugung hindeutet. Es besteht eine latente Spannung zwischen erlebter Kompetenzfrustrierung (zu niedriges KI-Level) und gleichzeitiger eigenständiger Anwendung — eine Mehrfachkodierung mit HK1 wäre denkbar, da die Person eigeninitiativ handelt, bleibt aber konservativ ausgeschlossen, da der Fokus klar auf Kompetenz liegt.}
\end{description}
\medskip

\begin{description}
  \item[\texttt{\#}\textbf{00:12:58}\texttt{\#} (Sinneinheit 15)]
  \begin{quote}
  Und da habe ich wirklich gesehen, dass Spalten teilweise komplett, obwohl sie in beiden Listen drinnen waren, weggelassen worden sind, somit dann die Gesamtergebnisse nicht mehr gestimmt haben. Und wenn ich das nicht im Detail alles wirklich quasi fast mit einem Leuchtstift durchkontrolliert hätte gedanklich, wäre ich komplett über dieses Ergebnis und über diese Auswertungen und Empfehlungen gestolpert. Und das ist für mich der größte Knackpunkt. Ich traue mir ehrlich gesagt bei solchen Themen, wenn es um detaillierte Auswertungen, KI-gestützte geht, nicht über die KI drüber. Ich habe manch andere Themen schon probiert. So simpler Büroalltag bei mir ist es, dass ich alle drei Monate, einmal im Quartal, so ein Minigolf-Turnier für unser Haus organisiere. Das tue ich für den Betriebsrat. Und da kriege ich diese Anmeldungen alle per Automaten-Mail.
  \end{quote}
  \textbf{Kategorie:} HK2 -- Kompetenzerleben\\
  \textbf{Weitere Kategorien:} HK1\\
  \textbf{Begr\"undung:} Die Passage beschreibt konkrete Fehler der KI bei detaillierten Auswertungen (fehlende Spalten, falsche Gesamtergebnisse), die die Interviewperson nur durch intensive Eigenkontrolle entdeckt. Dies verweist auf eine kritische Einschätzung der KI-Kompetenz und zugleich auf die Notwendigkeit eigener fachlicher Expertise, um KI-Outputs validieren zu können.
  \\\textit{Memo: Die Textstelle zeigt eine ambivalente Haltung: Einerseits wird die KI für bestimmte Aufgaben (komplexe Auswertungen) als unzuverlässig und fehleranfällig erlebt, andererseits wird ihr Einsatz für einfachere Routineaufgaben (Minigolf-Turnier, Büroalltag) angedeutet. Die Unterscheidung zwischen Aufgabentypen könnte in späteren Sinneinheiten als differenzierende Nutzungsstrategie weiter relevant werden.}
\end{description}
\medskip

\begin{description}
  \item[\texttt{\#}\textbf{00:12:58}\texttt{\#} (Sinneinheit 16)]
  \begin{quote}
  Wenn jemand ein Formular befüllt, kriege ich dazu das E-Mail mit der Anmeldung. Und da habe ich mir gedacht, Copilot-Routine kann ja nicht schwer sein, bastle mir aus all diesen E-Mails einfach eine Excel-Anmeldeliste. Und jedes Mal, wenn ich auf Knopfdruck, weil ein neues E-Mail da ist, schaut diese Anmeldeliste anders aus. Das ist jetzt aber auch kein Fehler der KI an dieser Stelle, möchte ich sagen, sondern ist mit Sicherheit ein Bedienfehler, weil ich das einfach zu schlecht geprompted habe. Zusammengefasst, alles, was in die detailliertere Analyse geht oder wo daraus wirklich konkrete Handlungen ausgelöst werden sollen, bin ich noch immer KI-skeptisch, auf Basis meiner mangelnden Kompetenz KI zu bedienen. Bei so Recherchethemen, wo es darum geht, neue Ideen zu entwickeln, Strategien einmal anzudenken, da liebe ich mittlerweile KI.
  \end{quote}
  \textbf{Kategorie:} HK2 -- Kompetenzerleben\\
  \textbf{Weitere Kategorien:} HK1\\
  \textbf{Begr\"undung:} Die Interviewperson beschreibt explizit ihre 'mangelnde Kompetenz KI zu bedienen' als Ursache für Misserfolge beim Prompting und zeigt damit eine kompetenzfrustrierende Erfahrung. Gleichzeitig wird eine kompetenzförderliche Dimension sichtbar, wenn KI für Recherche und Strategieentwicklung als wirksames Werkzeug erlebt wird.
  \\\textit{Memo: Die Textstelle zeigt eine differenzierte Selbsteinschätzung: Die Person unterscheidet klar zwischen Anwendungsbereichen, in denen KI als kompetenzförderlich erlebt wird (Ideenentwicklung, Recherche), und solchen, in denen Kompetenzdefizite im Umgang mit KI zu Frustration führen. Die Selbstattribution des Fehlers ('Bedienfehler') ist motivational relevant, da sie Kontrollüberzeugungen und Kompetenzwahrnehmung berührt.}
\end{description}
\medskip

\begin{description}
  \item[\texttt{\#}\textbf{00:17:48}\texttt{\#} (Sinneinheit 19)]
  \begin{quote}
  Ja, ich meine, wir haben ja bis vor ein, zwei Jahren ja alles ohne KI machen dürfen und müssen. Mir kommt vor, dass es früher teilweise mühsamer war. Und ich merke jetzt den Unterschied zwischen den Teammitgliedern, die mit KI arbeiten und den Mitgliedern, die weiterhin ohne KI arbeiten, aus welchen Gründen auch immer. Um mal schnell in die Gänge bei einem neuen Thema zu kommen, und bei mir sind oft eher neuere Themen am Tisch, ist KI mittlerweile für mich essentiell. Und ich gehe mittlerweile schon so weit, dass ich alle meine Kolleginnen und Kollegen dazu anhalte, bewusst mit KI zu arbeiten, damit sie sich selber einerseits Arbeitszeit ersparen, also effizienter werden, aber andererseits auch neue Sichtweisen gewinnen im Zusammenspiel mit der KI.
  \end{quote}
  \textbf{Kategorie:} HK2 -- Kompetenzerleben\\
  \textbf{Weitere Kategorien:} HK1, HK3\\
  \textbf{Begr\"undung:} Die Passage beschreibt explizit, wie KI bei neuen Themen als essenzielles Werkzeug erlebt wird, das Effizienz steigert und neue Sichtweisen eröffnet – beides verweist auf kompetenzförderliche Erfahrungen im Sinne erweiterter Wirksamkeit und verbesserter Ergebnisqualität.
  \\\textit{Memo: Die Sinneinheit ist thematisch dicht und verbindet mehrere Dimensionen: eigene Effizienzgewinne (HK2), selbstbestimmte Nutzungsentscheidung und Multiplikatorrolle (HK1), sowie Teamdifferenzierung und Führungsimpuls (HK3). Keine HK0-Zuweisung notwendig, da alle relevanten Aspekte spezifischeren Kategorien zugeordnet werden können.}
\end{description}
\medskip

\begin{description}
  \item[\texttt{\#}\textbf{00:17:48}\texttt{\#} (Sinneinheit 20)]
  \begin{quote}
  Also für mich ist die KI immer so ein Thema, einerseits für die Recherche, also diese generative AI ist für mich so wichtig, aber natürlich auch diesen Automatisierungsanteil, den KI liefern kann. Also diese Mitarbeit im Team ist mir sehr, sehr wichtig. Und von dem her gesehen, früher hat es oft länger gedauert, kommt mir vor, bis man erste Ergebnisse für Besprechungen gehabt hat. Heutzutage kann ich eigentlich schon relativ schnell mit ersten Inputs einmal arbeiten. Und auch in der Weiterverarbeitung früher, nur ein Beispiel, wenn wir unsere Vorstände zum Beispiel zu den Bundesgremien entsendet haben, haben wir früher die Präsentationen im Haus an alle Fachbereiche verteilt und alle zum Beispiel acht Fachbereiche mussten zurückmelden, was in dieser Präsentation drinnen ist und welche Sichtweisen sie haben.
  \end{quote}
  \textbf{Kategorie:} HK2 -- Kompetenzerleben\\
  \textbf{Weitere Kategorien:} HK0\\
  \textbf{Begr\"undung:} Die Interviewperson beschreibt, dass sie durch den Einsatz generativer KI deutlich schneller erste Ergebnisse und Inputs für Besprechungen erzielen kann als früher. Dies verweist auf eine verbesserte Wirksamkeitserfahrung und eine kompetenzförderliche Erweiterung der Handlungsmöglichkeiten durch KI.
  \\\textit{Memo: Die Sinneinheit bricht ab, bevor das konkrete Beispiel (Vorstände, Bundesgremien, Präsentationen) abgeschlossen wird. Der Vergleich 'früher–heute' ist ein klares Indiz für erlebten Kompetenz- und Effizienzgewinn. Die Textstelle enthält noch keinen expliziten Bezug zu Autonomie oder sozialer Eingebundenheit, obwohl der Verweis auf 'Mitarbeit im Team' als wichtig einen möglichen HK3-Anknüpfungspunkt für Folgepassagen andeutet.}
\end{description}
\medskip

\begin{description}
  \item[\texttt{\#}\textbf{00:17:48}\texttt{\#} (Sinneinheit 21)]
  \begin{quote}
  Mittlerweile machen wir zumindest diese erste Zusammenfassung, das erste Briefing mit Copilot und das auf SharePoint-Basis, damit es im gesicherten Bereich stattfindet. Aber du kannst den Vorstand viel besser vorbereiten, weil du dort acht bis zehn Präsentationen mit X-Folien gibst, einen Two-Pager mit, wo einmal eine Summary der Präsentation dabei ist und zusätzlich dann das Ergebnis der Fachbereichskollegen, das Feedback aus dem Fachbereich in Ergänzung. Und der Vorstand kann sich wirklich komprimiert einfach schnell vorbereiten. Das war früher mühsam. Und so, sage ich mal, ich könnte mir ehrlich gesagt die Arbeit ohne KI gar nicht mehr vorstellen. Wichtig ist mir halt, dass die Entscheidungsfindung trotzdem aber nur rein auf Basis valider Fakten und Einschätzungen von den Menschen getroffen werden.
  \end{quote}
  \textbf{Kategorie:} HK2 -- Kompetenzerleben\\
  \textbf{Weitere Kategorien:} HK1, HK0\\
  \textbf{Begr\"undung:} Die Aussage 'du kannst den Vorstand viel besser vorbereiten' und 'ich könnte mir die Arbeit ohne KI gar nicht mehr vorstellen' verweist auf eine wahrgenommene Erweiterung der eigenen Wirksamkeit und Ergebnisqualität durch den KI-Einsatz, was charakteristisch für kompetenzförderliche Erfahrungen ist.
  \\\textit{Memo: Die Sinneinheit enthält eine bemerkenswerte Spannungsdynamik: einerseits starke Abhängigkeit von und Begeisterung für KI ('könnte mir die Arbeit ohne KI gar nicht mehr vorstellen'), andererseits eine bewusste Grenzziehung hinsichtlich menschlicher Entscheidungshoheit. Diese Spannung zwischen KI-Nutzung und Erhalt humaner Kontrolle könnte für die Analyse des Autonomieerlebens im weiteren Verlauf besonders relevant sein.}
\end{description}
\medskip

\begin{description}
  \item[\texttt{\#}\textbf{00:17:48}\texttt{\#} (Sinneinheit 22)]
  \begin{quote}
  Also das, was ich zum Beispiel überhaupt nicht schätzen würde, wenn eine Kollegin, ein Kollege zu mir kommt mit einem Ergebnis, mir verschweigt, dass es KI generiert ist und dann daraus noch sogar eine Handlungsempfehlung ableitet, die er selber nicht durchdacht hat. Also das würde bei mir zu Unmut führen.
  \end{quote}
  \textbf{Kategorie:} HK2 -- Kompetenzerleben\\
  \textbf{Weitere Kategorien:} HK3\\
  \textbf{Begr\"undung:} Die Interviewperson kritisiert, dass Handlungsempfehlungen unreflektiert aus KI-Outputs übernommen werden, ohne eigenständiges Durchdenken. Dies verweist auf eine Sorge um Kompetenzattribution und professionelle Eigenleistung: Wenn KI-Ergebnisse als eigene Expertise ausgegeben werden, wird die tatsächliche fachliche Kompetenz der Beteiligten verschleiert und damit die Grundlage für fundierte Kompetenzbewertung untergraben.
  \\\textit{Memo: Die Sinneinheit berührt gleichzeitig die Dimension interpersonalen Vertrauens (HK3) und die Frage nach authentischer Kompetenzdemonstration (HK2). Ein möglicher HK1-Bezug (Kontrollverlust über Entscheidungsgrundlagen) ist denkbar, aber zu indirekt für eine sichere Kodierung. Der Begriff 'Unmut' signalisiert emotionale Frustration als Reaktion auf die beschriebene Situation.}
\end{description}
\medskip

\begin{description}
  \item[\texttt{\#}\textbf{00:18:49}\texttt{\#} (Sinneinheit 23)]
  \begin{quote}
  Nein, ich sage mal, das müsste durch seine eigenen Gedanken und Erfahrungen vor allem, die die KI ja gar nicht in unserem Haus haben kann, bei der Beurteilung von Sachverhalten ergänzt werden. Also für mich, ich weiß nicht, ob es der richtige Begriff ist, müsste jedes Ergebnis ein Hybrid-Ergebnis sein. Mensch und Maschine. Also rein mit einem maschinellen Ergebnis zu mir zu kommen, da muss ich mich echt hinterfragen, für was ich den Mitarbeiter noch brauche, weil das hätte ich dann selber auch machen können. Ich möchte nicht haben, dass meine Leute mit der KI, mit dem KI-Ergebnis in den Diskurs gehen, für sich selber. Und vielleicht mit ein, zwei anderen Kollegen auch noch mal kurz abprüfen, evaluieren. Und das Ergebnis könnten sie mir noch in kürzerer Zeit in einer höheren Wertigkeit, glaube ich, dann überliefern.
  \end{quote}
  \textbf{Kategorie:} HK2 -- Kompetenzerleben\\
  \textbf{Weitere Kategorien:} HK1, HK3\\
  \textbf{Begr\"undung:} Die Interviewperson betont explizit, dass ein rein maschinelles Ergebnis die professionelle Expertise der Mitarbeiter entwertet ('da muss ich mich echt hinterfragen, für was ich den Mitarbeiter noch brauche'). Gleichzeitig wird das Hybridmodell als Weg gesehen, durch menschliche Ergänzung eine höhere Ergebnisqualität zu erzielen – dies verweist auf Kompetenzerleben als Wirksamkeit und Expertise.
  \\\textit{Memo: Zentrale Passage zum Hybrid-Konzept: Die Führungskraft formuliert ein normatives Leitbild für den KI-Einsatz, das alle drei Grundbedürfnisse berührt. Besonders auffällig ist die implizite Sorge um Rollenentleerung und Expertiseentwertung (HK2), die eng mit der Frage nach dem Sinn menschlicher Mitarbeit verknüpft ist.}
\end{description}
\medskip

\begin{description}
  \item[\texttt{\#}\textbf{00:20:30}\texttt{\#} (Sinneinheit 24)]
  \begin{quote}
  Ja, dadurch, dass eigentlich dieser Arbeitsstil mittlerweile sehr breitflächig vertreten ist, wird wahrscheinlich von einigen in den Besprechungen schon fast als inflationär wahrgenommen. Vielleicht schon ein bisschen dahingehend interpretiert. Ja, da steckt eh keine persönliche Überlegung mehr dahinter. Das ist ein kurzer Prompt gewesen in ChatGPT, weil die Leute wissen ja nicht, ob ich es vielleicht mit einem Prompt in zwei Zeilen gelöst habe oder mit fünf Iterationen über sieben Tage verteilt, weil ich immer wieder das Ergebnis versucht habe zu verfeinern. Und ich glaube, wenn du mit einer KI-Präsentation aktuell auftrittst oder auch künftig, sollte man zu Beginn erklären, wie ist es zu diesem KI-generierten Ergebnis gekommen?
  \end{quote}
  \textbf{Kategorie:} HK2 -- Kompetenzerleben\\
  \textbf{Weitere Kategorien:} HK3\\
  \textbf{Begr\"undung:} Die Passage beschreibt eine Attributionsverschiebung: Die Interviewperson erlebt, dass ihre tatsächliche kognitive Leistung (mehrere Iterationen über sieben Tage) von anderen nicht wahrgenommen oder anerkannt wird, weil das Ergebnis pauschal als simpler KI-Output abgewertet wird. Dies verweist auf eine kompetenzfrustrierende Erfahrung, bei der die eigene Expertise und der investierte Aufwand unsichtbar werden.
  \\\textit{Memo: Die Sinneinheit berührt eine interessante Schnittfläche zwischen Kompetenz- und Eingebundenheitserleben: Die kompetenzfrustrierende Attributionsverschiebung entsteht primär durch den sozialen Bewertungskontext. Der Hinweis auf erklärende Transparenz ('wie ist es zu diesem KI-generierten Ergebnis gekommen') kann als Bewältigungsstrategie für beide Bedürfnisbereiche gelesen werden.}
\end{description}
\medskip

\begin{description}
  \item[\texttt{\#}\textbf{00:20:30}\texttt{\#} (Sinneinheit 25)]
  \begin{quote}
  Also wie gesagt, ich weiß jetzt nicht im Detail, welche Prompts ich verwendet habe, aber einfach, welche Anläufe ich genommen habe und auch zu symbolisieren, dass das jetzt nicht eine Arbeit von einer halben Stunde war, sondern dass es in Summe über eine Woche verteilt mehrere Iterationen waren und auch die Quellen sehr valide waren nach der Prüfung. Und dementsprechend muss man die Wertigkeit von einem KI-Ergebnis wahrscheinlich besser darstellen, sonst wird es als 0815-Prompt abgetan.
  \end{quote}
  \textbf{Kategorie:} HK2 -- Kompetenzerleben\\
  \textbf{Weitere Kategorien:} HK3\\
  \textbf{Begr\"undung:} Die Interviewperson beschreibt den iterativen Aufwand und die Qualitätsprüfung, die hinter einem KI-Ergebnis stehen – dies verweist auf das Erleben von Expertise und professionellem Einsatz. Gleichzeitig thematisiert sie die fehlende Anerkennung dieser Kompetenzleistung, was auf eine Frustration durch Attributionsverschiebung (Leistung wird dem Tool, nicht der Person zugeschrieben) hindeutet.
  \\\textit{Memo: Die Textstelle zeigt eine charakteristische Spannung zwischen tatsächlichem Kompetenzaufwand und dessen sozialer Wahrnehmung. Die Mehrfachkodierung HK2+HK3 ist gerechtfertigt, da die Kompetenzfrustration und die Anerkennungsproblematik analytisch trennbar, aber inhaltlich eng verknüpft sind. Relevant für die Frage, wie KI-Arbeit in Organisationen sichtbar gemacht und bewertet wird.}
\end{description}
\medskip

\begin{description}
  \item[\texttt{\#}\textbf{00:21:15}\texttt{\#} (Sinneinheit 26)]
  \begin{quote}
  Man muss es sogar wirklich am Anfang dazu sagen, dass man einerseits seine eigene Kompetenz in der Anwendung von KI dem Auditorium näherbringt, dass die nicht glauben, das war einfach beim Autofahren über das Smartphone ein diktierter Prompt, ein zweizeiliger. Sondern wie gesagt, man hat das methodisch kompetent vorgenommen, dieses KI-generierte Ergebnis zu gewinnen. Weil sonst, glaube ich, zweifeln viele schon bei der Ergebnispräsentation daran, weil sie sich denken, das ist ein 0815-Prompt und die Quellen waren sicher nicht valide.
  \end{quote}
  \textbf{Kategorie:} HK2 -- Kompetenzerleben\\
  \textbf{Weitere Kategorien:} HK3\\
  \textbf{Begr\"undung:} Die Passage beschreibt, wie die Interviewperson aktiv ihre methodische Kompetenz im Umgang mit KI gegenüber dem Auditorium kommunizieren muss, um Glaubwürdigkeit zu sichern. Dies verweist auf das Bedürfnis, als kompetenter Anwender wahrgenommen zu werden und einer drohenden Expertiseentwertung durch oberflächliche KI-Nutzungswahrnehmung entgegenzuwirken.
  \\\textit{Memo: Die Passage zeigt eine interessante Verschränkung von Kompetenz- und Eingebundenheitserleben: Die Notwendigkeit, KI-Kompetenz explizit zu kommunizieren, entsteht aus einem sozialen Legitimitätsproblem. Der Rechtfertigungsdruck könnte auch als leicht autonomiefrustrierend interpretiert werden, wird hier jedoch konservativ nicht als HK1 codiert, da der Fokus klar auf Kompetenzdarstellung und sozialer Wahrnehmung liegt.}
\end{description}
\medskip

\begin{description}
  \item[\texttt{\#}\textbf{00:23:20}\texttt{\#} (Sinneinheit 28)]
  \begin{quote}
  Weil wenn ich dann wirklich Kolleginnen und Kollegen mit teilweise 20-jährigem Know-how, sprich Wissen und Erfahrung habe, plus das ganze Agenze durch schlaue Agenten, die einfach viele so bürokratische oder vielleicht auch Recherchetätigkeiten im Unternehmenswissen, von der Rechtsabteilung bis zum Datenschutz bis zum Marketing für uns erledigen können, dann kann er viel schlagkräftiger werden in unserem Haus. Und es gibt zum Beispiel für mich Beispiele, wo ich mit menschlichem Aufwand viel zu viele Kosten in die Hand nehmen würde, wenn ich das einmal schlau über AI-Netzwerke, über kleine mit Automatisierung löse, kann ich auf einmal zum Beispiel einen Newsletter für steirische Firmen und Kundenberater zu den wesentlichsten Wirtschafts-News der Steiermark kostengünstigst erwirken, ganzjährig.
  \end{quote}
  \textbf{Kategorie:} HK2 -- Kompetenzerleben\\
  \textbf{Weitere Kategorien:} HK1\\
  \textbf{Begr\"undung:} Die Passage beschreibt eine deutliche Erweiterung der eigenen Wirksamkeit: Die Kombination aus menschlichem 20-jährigen Know-how und KI-Agenten ermöglicht eine erhöhte 'Schlagkraft', was auf kompetenzförderliche Erfahrungen durch erweiterte Fähigkeiten und verbesserte Ergebnisqualität hinweist.
  \\\textit{Memo: Die Textstelle verbindet instrumentelle Effizienzgewinne mit einem klaren Erleben erweiterter Handlungs- und Wirksamkeitspotenziale. Die Metapher 'schlagkräftiger werden' ist besonders aufschlussreich für das Kompetenzerleben. Eine Zuordnung zu HK3 wäre denkbar (Teamsynergie Mensch-KI), erscheint jedoch nicht hinreichend explizit belegt.}
\end{description}
\medskip

\begin{description}
  \item[\texttt{\#}\textbf{00:23:20}\texttt{\#} (Sinneinheit 29)]
  \begin{quote}
  Früher, wenn ich das gemacht hätte über Manpower oder Womenpower in meinem Team, da wäre es mir zu schade gewesen, diese wertvollen Ressourcen für so einen Newsletter einzusetzen.
  \end{quote}
  \textbf{Kategorie:} HK2 -- Kompetenzerleben\\
  \textbf{Weitere Kategorien:} HK1\\
  \textbf{Begr\"undung:} Die Formulierung 'wertvollen Ressourcen' impliziert eine professionelle Einschätzung, dass Teamkapazitäten für höherwertige Aufgaben reserviert werden sollten – dies deutet auf ein Kompetenzerleben hin, das mit der qualitativen Bewertung von Arbeitsinhalten und Wirksamkeit verbunden ist.
  \\\textit{Memo: Die Sinneinheit gewinnt ihren eigentlichen motivationalen Gehalt erst im Kontrast zur KI-gestützten Gegenwart: Implizit wird eine frühere Einschränkung (kein Newsletter wegen Ressourcenknappheit) beschrieben, die durch KI-Einsatz aufgehoben wurde. Für die Vollständigkeit der Interpretation sollte der unmittelbare Kontext (Folgesätze) herangezogen werden.}
\end{description}
\medskip

\begin{description}
  \item[\texttt{\#}\textbf{00:24:20}\texttt{\#} (Sinneinheit 30)]
  \begin{quote}
  Und ich kann was verdienen damit. Vielleicht noch eine Rechnung. Wenn ich so einen Newsletter, das ist ein reiner Proof-of-Concept-Intern für steirische Firmen und Kundenberater, wo einfach die wesentlichsten steirischen, aber auch österreichischen Medien zu Wirtschaftsthemen abgegriffen werden und daraus branchenspezifische News einmal in der Woche, am Montag in der Früh ausgespielt werden, sowohl als Newsletter schriftlich, als aber auch als Podcast. Das kostet mir im ersten Schritt keine 14.000 Euro über externe Programmierung außerhalb unserer Systeme. Dann brauche ich noch das Hosting, das kostet 1.000 Euro, da bin ich bei 15.000 brutto.
  \end{quote}
  \textbf{Kategorie:} HK2 -- Kompetenzerleben\\
  \textbf{Weitere Kategorien:} HK0\\
  \textbf{Begr\"undung:} Die detaillierte Kostenkalkulation und die eigenständige Konzeption des KI-gestützten Projekts (Newsletter + Podcast aus Medienquellen) deuten auf ein Wirksamkeitserleben hin: Die Interviewperson präsentiert sich als kompetente Gestalterin, die durch KI-Einsatz eine aufwendige externe Lösung kostengünstig realisieren kann.
  \\\textit{Memo: Die einleitende Aussage 'Ich kann was verdienen damit' könnte auf Autonomie- oder Kompetenzerleben hinweisen, bleibt aber zu vage für eine sichere HK1-Zuweisung. Der Hauptfokus liegt auf der technisch-wirtschaftlichen Projektbeschreibung. Möglicher Querverweis auf HK1, falls in benachbarten Sinneinheiten Eigeninitiative oder Entscheidungsfreiheit bei der Projektumsetzung thematisiert wird.}
\end{description}
\medskip

\begin{description}
  \item[\texttt{\#}\textbf{00:24:20}\texttt{\#} (Sinneinheit 31)]
  \begin{quote}
  Und wenn ich es schaffe, dass jede steirische Raiffeisenbank in weiterer Folge von dieser Dienstleistung profitieren möchte und 50 Euro im Monat, egal wie viele Users sie hat, bezahlt, bin ich bei 40.000, 50.000 und 2.000 Euro, bin ich bei 24.000 Euro Erträgen, obwohl ich einmal im ersten Schritt circa 15.000 investiert habe und das im ersten Jahr. Und das sind Themen, die würde ich mit menschlichem Arbeitseinsatz nie profitabel gestalten können.
  \end{quote}
  \textbf{Kategorie:} HK2 -- Kompetenzerleben\\
  \textbf{Weitere Kategorien:} HK1\\
  \textbf{Begr\"undung:} Die Passage beschreibt, wie der Einsatz generativer KI die Ergebnisqualität und Wirksamkeit des eigenen Handelns deutlich steigert – ein mit menschlichem Arbeitseinsatz allein nicht erreichbares Ergebnis wird durch KI realisierbar. Dies verweist auf eine kompetenzförderliche Erfahrung im Sinne erweiterter Handlungsfähigkeit und verbesserter Outcomes.
  \\\textit{Memo: Die Textstelle enthält eine konkrete Kalkulation als Beleg für den wahrgenommenen KI-Mehrwert. Die Verbindung von Kompetenz- und Autonomieerleben ist hier eng verknüpft: Die KI erweitert nicht nur die Fähigkeiten, sondern eröffnet auch einen selbstgestalteten unternehmerischen Spielraum. Eine Zuordnung zu HK0 (institutioneller Kontext Raiffeisenbanken) wäre denkbar, tritt jedoch hinter die motivationale Dimension zurück.}
\end{description}
\medskip

\begin{description}
  \item[\texttt{\#}\textbf{00:26:32}\texttt{\#} (Sinneinheit 33)]
  \begin{quote}
  Ja, also mir ist wichtig, dass wir vor allem bestmöglich auch diese KI-Möglichkeiten gezielt einsetzen, eben von der Generierung bis zur Automatisierung und so weiter, aber trotzdem permanent Messpunkte entwickeln, um zu schauen, ob wir dadurch wirklich effizienter und effektiver werden. Und vor allem, wenn wirklich losgelöste, autonome Agenten dann für uns Tätigkeiten machen, wie können wir wirklich permanent bewerten, ob das, was diese Zusatzkollegen für uns generieren und produzieren, wirklich in das, was wir erreichen wollen, die strategischen Ziele, wirklich auch einzahlt. Wenn es nur operative Entlastung ist, ist es auch gut, aber es muss eine gewisse Gratwanderung sein, bis wohin hilft es uns wirklich weiter und wo ist vielleicht zu viel des Guten.
  \end{quote}
  \textbf{Kategorie:} HK2 -- Kompetenzerleben\\
  \textbf{Weitere Kategorien:} HK1\\
  \textbf{Begr\"undung:} Die Passage thematisiert die Notwendigkeit, die Qualität und Wirksamkeit KI-generierter Outputs permanent zu bewerten und deren Beitrag zu strategischen Zielen zu messen, was auf ein professionelles Kompetenzverständnis verweist. Die Reflexion über 'zu viel des Guten' zeigt ein Bewusstsein für die Grenzen der KI-Nutzung und die eigene Rolle als kompetenter Beurteilungsinstanz.
  \\\textit{Memo: Die Sinneinheit enthält eine charakteristische Ambivalenz: KI-Nutzung wird begrüßt (operative Entlastung), gleichzeitig wird die Notwendigkeit eigenständiger Steuerungs- und Bewertungskompetenz betont. Die Formulierung 'Zusatzkollegen' ist semantisch auffällig und könnte in Bezug auf HK3 (soziale Eingebundenheit, Mensch-KI-Beziehung) in späteren Durchgängen weiter reflektiert werden.}
\end{description}
\medskip

\begin{description}
  \item[\texttt{\#}\textbf{00:26:32}\texttt{\#} (Sinneinheit 34)]
  \begin{quote}
  Ich sage nur zum Beispiel das Thema Tokenisierung oder Coins, wenn das irgendwann einmal logischerweise mit monetären Kosten, mit einem monetären Aufwand verbunden ist, dann kann es uns auch passieren, dass das, was der Agent für uns macht, in Wahrheit teurer ist, weil es ein 0815-Thema ist, mit schlechter Qualität, als wenn es der Mensch weiterhin für uns gemacht hätte. Hört sich zwar ein bisschen perfide an, aber die nächste Herausforderung wird dann werden, ist dieser Agenteneinsatz wirklich effektiv und ist er effizient und wäre es nicht besser gewesen, das Thema wieder abzudrehen oder durch einen Menschen zu machen?
  \end{quote}
  \textbf{Kategorie:} HK2 -- Kompetenzerleben\\
  \textbf{Weitere Kategorien:} HK0, HK1\\
  \textbf{Begr\"undung:} Die Thematisierung der Qualitätseinschätzung ('schlechte Qualität', 'wäre es nicht besser gewesen') deutet auf eine kompetenzrelevante Abwägung hin: Die Interviewperson reflektiert kritisch, ob KI-generierte Ergebnisse menschlicher Expertise überlegen sind, was eine Auseinandersetzung mit möglicher Expertiseentwertung oder -bestätigung impliziert.
  \\\textit{Memo: Die Textstelle enthält eine nuancierte Kosten-Nutzen-Reflexion über KI-Agenten, die gleichzeitig Kontextfaktoren (Effizienz, Tokenisierungskosten), Autonomie (Entscheidung über Ab- oder Weiterbetrieb) und Kompetenz (Qualitätsvergleich Mensch vs. KI) berührt. Die Mehrfachkodierung erscheint gerechtfertigt, da alle drei Dimensionen substanziell angesprochen werden.}
\end{description}
\medskip

% ── HK3: Soziale Eingebundenheit ──────────────────────────────────────────────
\subsection*{HK3 -- Soziale Eingebundenheit}

\begin{description}
  \item[\texttt{\#}\textbf{00:04:52}\texttt{\#} (Sinneinheit 6)]
  \begin{quote}
  Dementsprechend habe ich dann selber sehr viel mit der ganzen M365-Kollaborationssoftware zu tun. Und was ich immer früher gehabt habe in meinem Aufgabenbereich, brauche ich sehr oft Präsentationen. Präsentationen vor allem auch zu Themen, die es bei uns im Haus noch nicht gibt. Und ich bin immer an, ich versuche möglichst weit nach vorne zu denken. Und früher hätte ich da immer zwei, drei meiner Leute beauftragt, mir einmal das Thema zu recherchieren. Heutzutage kann ich da ganz offen und ehrlich sagen, mache ich das nur mit KI, mit klassisch, wenn es etwas ist, was ich außerhalb der Unternehmenssecurity machen darf, weil es ein ganz neues Thema ist, wie zum Beispiel Lebensphasenkonzepte im Vertrieb.
  \end{quote}
  \textbf{Kategorie:} HK3 -- Soziale Eingebundenheit\\
  \textbf{Weitere Kategorien:} HK1, HK2\\
  \textbf{Begr\"undung:} Implizit beschreibt die Interviewperson eine Veränderung der Zusammenarbeit: Aufgaben, die früher gemeinsam mit Teammitgliedern erledigt wurden, übernimmt nun die KI allein. Dies deutet auf eine Verschiebung in der Teaminteraktion und potenziell veränderte Rollenbeziehungen hin.
  \\\textit{Memo: Die Textstelle enthält einen relevanten Nebenbefund: Die explizite Erwähnung der 'Unternehmenssecurity' als Nutzungsbeschränkung könnte auch HK0 (regulatorisch-organisationaler Kontext) rechtfertigen, wurde jedoch hier nicht separat vergeben, da sie nur als Nebenbedingung erwähnt wird und der Fokus der Passage auf dem veränderten Arbeitsverhalten liegt. HK3 ist schwach ausgeprägt und könnte bei Revision als zu spekulativ eingestuft werden.}
\end{description}
\medskip

\begin{description}
  \item[\texttt{\#}\textbf{00:07:54}\texttt{\#} (Sinneinheit 9)]
  \begin{quote}
  Ganz schräg, ich kann mittlerweile nicht mehr differenzieren, ob man gegenüber dieser KI, egal ob ich mit ihr spreche, mit ihr schreibe, ein Mensch ist oder kein Mensch ist. Was meine ich damit? Mir ist das ehrlich gesagt echt egal. Ich finde es, das ist ganz, für mich echt irritierend oft, ich bin ein sehr digitalaffiner Mensch und ich tue mir oft leichter, sogar mit der KI zu kommunizieren. Warum? Weil es ist oft so mühsam, einen Kollegen zum Zeitpunkt erst für einen Termin zu bekommen oder ihn jetzt anzurufen und er ist eh nicht erreichbar. Die KI ist für mich immer da, das ist für mich ein Convenience-Thema und ich kriege teilweise von einer KI Feedbacks oder Meinungen zugespiegelt, da muss ich vorher natürlich achten, wo kommt die Recherche her.
  \end{quote}
  \textbf{Kategorie:} HK3 -- Soziale Eingebundenheit\\
  \textbf{Weitere Kategorien:} HK1\\
  \textbf{Begr\"undung:} Die Interviewperson beschreibt explizit eine veränderte Beziehungsdynamik: Sie kommuniziert leichter mit der KI als mit Kollegen und empfindet die zwischenmenschliche Koordination als mühsam. Dies verweist auf eine Verschiebung in der sozialen Eingebundenheit, bei der die KI zunehmend interpersonelle Interaktionen ersetzt oder verdrängt.
  \\\textit{Memo: Interessante Ambiguität: Die erlebte 'Erleichterung' im Umgang mit KI gegenüber Kollegen könnte sowohl als Autonomiegewinn (HK1) als auch als möglicher Rückzug aus sozialen Beziehungen (HK3, potenziell frustrierend) gelesen werden. Die Passage ist motivational bedeutsam, da sie auf eine subjektive Neubewertung sozialer Interaktion hinweist. Die Aussage 'mir ist es egal, ob Mensch oder nicht' deutet auf eine bemerkenswerte Verschiebung in der wahrgenommenen Qualität von Beziehungen hin.}
\end{description}
\medskip

\begin{description}
  \item[\texttt{\#}\textbf{00:07:54}\texttt{\#} (Sinneinheit 10)]
  \begin{quote}
  Ich versuche natürlich immer einen Prompt so darzustellen, dass der Prompt nur mehr valide Quellen berücksichtigt und so weiter mit Quellenangaben. Aber für mich ist die KI mittlerweile bei vielen Themen ein gleichwertiger Arbeitskollege oder Arbeitspartner und hat den großen Charme, dass ich manchen physischen Kollegen nicht mehr brauche. Ich liebe Menschen. Aber bei meinen Themen ist es oft lustiger, mit der KI zu arbeiten.
  \end{quote}
  \textbf{Kategorie:} HK3 -- Soziale Eingebundenheit\\
  \textbf{Weitere Kategorien:} HK1, HK2\\
  \textbf{Begr\"undung:} Die Aussage, dass physische Kollegen teilweise durch die KI ersetzt werden, und die Charakterisierung der KI als 'gleichwertiger Arbeitskollege' beschreibt eine Verschiebung in der sozialen Interaktionsstruktur. Dies berührt direkt das Erleben von Zugehörigkeit und Beziehungsgestaltung im organisationalen Kontext.
  \\\textit{Memo: Bemerkenswerte Ambivalenz: Die Interviewperson betont 'Ich liebe Menschen', relativiert dies aber unmittelbar durch die Präferenz für KI-Interaktion. Die KI wird explizit als sozialer Akteur ('Arbeitskollege') konzeptualisiert – dies könnte auf ein ungewöhnlich hohes Maß an parasozialer Bindung an die KI hindeuten, das im weiteren Verlauf des Interviews zu beobachten wäre.}
\end{description}
\medskip

\begin{description}
  \item[\texttt{\#}\textbf{00:14:36}\texttt{\#} (Sinneinheit 17)]
  \begin{quote}
  Also dadurch, dass ich momentan noch so plump auftrete in Form dessen, dass das KI-Ergebnis, das mehrfach geprompte und iterierte, dann im Notebook LM einfach in so eine sehr bunte PowerPoint oder bunte PDF, muss ich sagen, gegossen wird, sage ich das immer gleich zu Beginn bei meinen Präsentationen dazu. Das Thema, wo ich immer sehr schnell versuche, mit KI erst zu recherchieren, das Ergebnis ist ja bitte über Notebook LM entstanden, die Visualisierung. Also ich würde nie verschweigen, dass ich KI verwendet habe dazu. Das ist mir sogar ganz, ganz wichtig, auch wenn es vielleicht beim Auditorium ein wenig die Skepsis erhöht. Aber in den meisten Fällen, wo ich mit KI jetzt aufgetreten bin und gerade die letzten zwei Tage in Salzburg auch mit in Summe zwei Präsentationen, hat sich niemand dadurch gestört gefühlt, dass das KI-aufbereitet war.
  \end{quote}
  \textbf{Kategorie:} HK3 -- Soziale Eingebundenheit\\
  \textbf{Weitere Kategorien:} HK1\\
  \textbf{Begr\"undung:} Die Interviewperson reflektiert, wie das Auditorium auf den offengelegten KI-Einsatz reagiert, und beschreibt das Ausbleiben negativer Reaktionen als positive soziale Erfahrung. Die Sorge um Skepsis und die Erleichterung über Akzeptanz verweisen auf das Erleben von Zugehörigkeit und Wertschätzung im organisationalen bzw. öffentlichen Kontext.
  \\\textit{Memo: Die Transparenz-Norm, die die Interviewperson für sich selbst setzt ('Das ist mir sogar ganz, ganz wichtig'), könnte auch auf Kompetenzerleben (professionelle Integrität) hinweisen. Die Passage bleibt jedoch primär im Bereich sozialer Akzeptanz und autonomer Werteorientierung; HK2 wird konservativ nicht vergeben.}
\end{description}
\medskip

\begin{description}
  \item[\texttt{\#}\textbf{00:17:48}\texttt{\#} (Sinneinheit 19)]
  \begin{quote}
  Ja, ich meine, wir haben ja bis vor ein, zwei Jahren ja alles ohne KI machen dürfen und müssen. Mir kommt vor, dass es früher teilweise mühsamer war. Und ich merke jetzt den Unterschied zwischen den Teammitgliedern, die mit KI arbeiten und den Mitgliedern, die weiterhin ohne KI arbeiten, aus welchen Gründen auch immer. Um mal schnell in die Gänge bei einem neuen Thema zu kommen, und bei mir sind oft eher neuere Themen am Tisch, ist KI mittlerweile für mich essentiell. Und ich gehe mittlerweile schon so weit, dass ich alle meine Kolleginnen und Kollegen dazu anhalte, bewusst mit KI zu arbeiten, damit sie sich selber einerseits Arbeitszeit ersparen, also effizienter werden, aber andererseits auch neue Sichtweisen gewinnen im Zusammenspiel mit der KI.
  \end{quote}
  \textbf{Kategorie:} HK3 -- Soziale Eingebundenheit\\
  \textbf{Weitere Kategorien:} HK1, HK2\\
  \textbf{Begr\"undung:} Die Interviewperson thematisiert den Unterschied zwischen KI-nutzenden und nicht-nutzenden Teammitgliedern und ihre Rolle dabei, Kolleg:innen aktiv zur KI-Nutzung anzuhalten – dies berührt Teamdynamiken und die Führungsbeziehung im Kontext von KI-Einführung.
  \\\textit{Memo: Die Sinneinheit ist thematisch dicht und verbindet mehrere Dimensionen: eigene Effizienzgewinne (HK2), selbstbestimmte Nutzungsentscheidung und Multiplikatorrolle (HK1), sowie Teamdifferenzierung und Führungsimpuls (HK3). Keine HK0-Zuweisung notwendig, da alle relevanten Aspekte spezifischeren Kategorien zugeordnet werden können.}
\end{description}
\medskip

\begin{description}
  \item[\texttt{\#}\textbf{00:17:48}\texttt{\#} (Sinneinheit 22)]
  \begin{quote}
  Also das, was ich zum Beispiel überhaupt nicht schätzen würde, wenn eine Kollegin, ein Kollege zu mir kommt mit einem Ergebnis, mir verschweigt, dass es KI generiert ist und dann daraus noch sogar eine Handlungsempfehlung ableitet, die er selber nicht durchdacht hat. Also das würde bei mir zu Unmut führen.
  \end{quote}
  \textbf{Kategorie:} HK3 -- Soziale Eingebundenheit\\
  \textbf{Weitere Kategorien:} HK2\\
  \textbf{Begr\"undung:} Die Passage beschreibt eine konkrete Erwartungshaltung an die Qualität der interpersonalen Zusammenarbeit: Transparenz und intellektuelle Eigenleistung werden als Grundlage einer vertrauensvollen Führungsbeziehung betrachtet. Das Verschweigen von KI-Generierung wird als Verletzung von Beziehungsnormen erlebt, was auf eine Beeinträchtigung der sozialen Eingebundenheit und Wertschätzung im Teamkontext verweist.
  \\\textit{Memo: Die Sinneinheit berührt gleichzeitig die Dimension interpersonalen Vertrauens (HK3) und die Frage nach authentischer Kompetenzdemonstration (HK2). Ein möglicher HK1-Bezug (Kontrollverlust über Entscheidungsgrundlagen) ist denkbar, aber zu indirekt für eine sichere Kodierung. Der Begriff 'Unmut' signalisiert emotionale Frustration als Reaktion auf die beschriebene Situation.}
\end{description}
\medskip

\begin{description}
  \item[\texttt{\#}\textbf{00:18:49}\texttt{\#} (Sinneinheit 23)]
  \begin{quote}
  Nein, ich sage mal, das müsste durch seine eigenen Gedanken und Erfahrungen vor allem, die die KI ja gar nicht in unserem Haus haben kann, bei der Beurteilung von Sachverhalten ergänzt werden. Also für mich, ich weiß nicht, ob es der richtige Begriff ist, müsste jedes Ergebnis ein Hybrid-Ergebnis sein. Mensch und Maschine. Also rein mit einem maschinellen Ergebnis zu mir zu kommen, da muss ich mich echt hinterfragen, für was ich den Mitarbeiter noch brauche, weil das hätte ich dann selber auch machen können. Ich möchte nicht haben, dass meine Leute mit der KI, mit dem KI-Ergebnis in den Diskurs gehen, für sich selber. Und vielleicht mit ein, zwei anderen Kollegen auch noch mal kurz abprüfen, evaluieren. Und das Ergebnis könnten sie mir noch in kürzerer Zeit in einer höheren Wertigkeit, glaube ich, dann überliefern.
  \end{quote}
  \textbf{Kategorie:} HK3 -- Soziale Eingebundenheit\\
  \textbf{Weitere Kategorien:} HK1, HK2\\
  \textbf{Begr\"undung:} Die Passage thematisiert die Interaktion zwischen Mitarbeitern und der Führungskraft sowie kollegialen Austausch ('mit ein, zwei anderen Kollegen auch noch mal kurz abprüfen, evaluieren'). Die Erwartung an einen gemeinsamen Diskursprozess verweist auf die Bedeutung sozialer Eingebundenheit im Arbeitskontext mit KI.
  \\\textit{Memo: Zentrale Passage zum Hybrid-Konzept: Die Führungskraft formuliert ein normatives Leitbild für den KI-Einsatz, das alle drei Grundbedürfnisse berührt. Besonders auffällig ist die implizite Sorge um Rollenentleerung und Expertiseentwertung (HK2), die eng mit der Frage nach dem Sinn menschlicher Mitarbeit verknüpft ist.}
\end{description}
\medskip

\begin{description}
  \item[\texttt{\#}\textbf{00:20:30}\texttt{\#} (Sinneinheit 24)]
  \begin{quote}
  Ja, dadurch, dass eigentlich dieser Arbeitsstil mittlerweile sehr breitflächig vertreten ist, wird wahrscheinlich von einigen in den Besprechungen schon fast als inflationär wahrgenommen. Vielleicht schon ein bisschen dahingehend interpretiert. Ja, da steckt eh keine persönliche Überlegung mehr dahinter. Das ist ein kurzer Prompt gewesen in ChatGPT, weil die Leute wissen ja nicht, ob ich es vielleicht mit einem Prompt in zwei Zeilen gelöst habe oder mit fünf Iterationen über sieben Tage verteilt, weil ich immer wieder das Ergebnis versucht habe zu verfeinern. Und ich glaube, wenn du mit einer KI-Präsentation aktuell auftrittst oder auch künftig, sollte man zu Beginn erklären, wie ist es zu diesem KI-generierten Ergebnis gekommen?
  \end{quote}
  \textbf{Kategorie:} HK3 -- Soziale Eingebundenheit\\
  \textbf{Weitere Kategorien:} HK2\\
  \textbf{Begr\"undung:} Die Passage thematisiert den sozialen Wahrnehmungskontext in Besprechungen: Die Interviewperson fürchtet, dass KI-gestützte Ergebnisse als mangelnde persönliche Auseinandersetzung interpretiert werden und dadurch Wertschätzung und Anerkennung im Team ausbleiben. Die vorgeschlagene Transparenz über den Entstehungsprozess zielt darauf ab, diese soziale Eingebundenheit und Glaubwürdigkeit zu erhalten.
  \\\textit{Memo: Die Sinneinheit berührt eine interessante Schnittfläche zwischen Kompetenz- und Eingebundenheitserleben: Die kompetenzfrustrierende Attributionsverschiebung entsteht primär durch den sozialen Bewertungskontext. Der Hinweis auf erklärende Transparenz ('wie ist es zu diesem KI-generierten Ergebnis gekommen') kann als Bewältigungsstrategie für beide Bedürfnisbereiche gelesen werden.}
\end{description}
\medskip

\begin{description}
  \item[\texttt{\#}\textbf{00:20:30}\texttt{\#} (Sinneinheit 25)]
  \begin{quote}
  Also wie gesagt, ich weiß jetzt nicht im Detail, welche Prompts ich verwendet habe, aber einfach, welche Anläufe ich genommen habe und auch zu symbolisieren, dass das jetzt nicht eine Arbeit von einer halben Stunde war, sondern dass es in Summe über eine Woche verteilt mehrere Iterationen waren und auch die Quellen sehr valide waren nach der Prüfung. Und dementsprechend muss man die Wertigkeit von einem KI-Ergebnis wahrscheinlich besser darstellen, sonst wird es als 0815-Prompt abgetan.
  \end{quote}
  \textbf{Kategorie:} HK3 -- Soziale Eingebundenheit\\
  \textbf{Weitere Kategorien:} HK2\\
  \textbf{Begr\"undung:} Die Aussage, dass das KI-Ergebnis als '0815-Prompt abgetan' wird, beschreibt ein Erleben mangelnder organisationaler Wertschätzung gegenüber der geleisteten Arbeit. Dies berührt die soziale Eingebundenheit, da die Interviewperson eine fehlende Anerkennung im sozialen/organisationalen Kontext wahrnimmt.
  \\\textit{Memo: Die Textstelle zeigt eine charakteristische Spannung zwischen tatsächlichem Kompetenzaufwand und dessen sozialer Wahrnehmung. Die Mehrfachkodierung HK2+HK3 ist gerechtfertigt, da die Kompetenzfrustration und die Anerkennungsproblematik analytisch trennbar, aber inhaltlich eng verknüpft sind. Relevant für die Frage, wie KI-Arbeit in Organisationen sichtbar gemacht und bewertet wird.}
\end{description}
\medskip

\begin{description}
  \item[\texttt{\#}\textbf{00:21:15}\texttt{\#} (Sinneinheit 26)]
  \begin{quote}
  Man muss es sogar wirklich am Anfang dazu sagen, dass man einerseits seine eigene Kompetenz in der Anwendung von KI dem Auditorium näherbringt, dass die nicht glauben, das war einfach beim Autofahren über das Smartphone ein diktierter Prompt, ein zweizeiliger. Sondern wie gesagt, man hat das methodisch kompetent vorgenommen, dieses KI-generierte Ergebnis zu gewinnen. Weil sonst, glaube ich, zweifeln viele schon bei der Ergebnispräsentation daran, weil sie sich denken, das ist ein 0815-Prompt und die Quellen waren sicher nicht valide.
  \end{quote}
  \textbf{Kategorie:} HK3 -- Soziale Eingebundenheit\\
  \textbf{Weitere Kategorien:} HK2\\
  \textbf{Begr\"undung:} Die Textstelle thematisiert die soziale Dynamik bei der Ergebnispräsentation: Die Interviewperson schildert, dass das Auditorium die KI-Nutzung skeptisch bewertet und die Wertschätzung der Arbeit von der wahrgenommenen Methodik abhängt. Dies berührt das Erleben von Anerkennung und Zugehörigkeit im organisationalen Kontext.
  \\\textit{Memo: Die Passage zeigt eine interessante Verschränkung von Kompetenz- und Eingebundenheitserleben: Die Notwendigkeit, KI-Kompetenz explizit zu kommunizieren, entsteht aus einem sozialen Legitimitätsproblem. Der Rechtfertigungsdruck könnte auch als leicht autonomiefrustrierend interpretiert werden, wird hier jedoch konservativ nicht als HK1 codiert, da der Fokus klar auf Kompetenzdarstellung und sozialer Wahrnehmung liegt.}
\end{description}
\medskip

\begin{description}
  \item[\texttt{\#}\textbf{00:23:20}\texttt{\#} (Sinneinheit 27)]
  \begin{quote}
  Für mich ist KI mittlerweile ein wesentlicher Bestandteil unseres täglichen, beziehungsweise zumindest wöchentlichen Tuns. Jeder muss sich in seinem Rollenprofil auch überlegen, wo er KI wirklich fix auch einbauen möchte, um seine Aufgaben zu erledigen. Jeder sollte bei mir im Team ganz klar überlegen, wie kann er KI nutzen, um effizienter, aber auch effektiver oder innovativer zu werden. Dementsprechend ist für mich auch mittlerweile, wenn dieses Thema Agentic AI bei uns wirklich hoffentlich stärker kommen sollte in den nächsten Jahren, muss man wirklich sich überlegen, wie man die Teams hybrid zusammensetzt.
  \end{quote}
  \textbf{Kategorie:} HK3 -- Soziale Eingebundenheit\\
  \textbf{Weitere Kategorien:} HK0, HK1\\
  \textbf{Begr\"undung:} Der Verweis auf hybride Teamzusammensetzungen im Kontext von Agentic AI deutet auf eine antizipierte Veränderung der Teamstruktur und der sozialen Zusammenarbeit hin, was die soziale Eingebundenheit im organisationalen Kontext direkt berührt.
  \\\textit{Memo: Die Sinneinheit verbindet Kontextbeschreibung (KI als etablierter Bestandteil) mit einer autonomieorientierten Führungshaltung und einem Ausblick auf veränderte Teamdynamiken. Besonders auffällig ist die normative Erwartungshaltung ('jeder muss', 'jeder sollte'), die eine gewisse Spannung zwischen proklamierter Autonomieförderung und implizitem Anpassungsdruck erzeugt – relevant für spätere Feincodierung.}
\end{description}
\medskip

\begin{description}
  \item[\texttt{\#}\textbf{00:24:58}\texttt{\#} (Sinneinheit 32)]
  \begin{quote}
  Ich bin leider viel zu sehr Science-Fiction-Fan, aber für mich ist das irgendwie so wie bei Star Wars, wo Menschen mit Robotern ihr Leben teilen und wo das eigentlich wirklich über die rationale Beziehung hinaus teilweise wirklich in eine emotionale Verbindung geht. Ich bin so ein Typ, ich bin leider für das zu sehr geeignet.
  \end{quote}
  \textbf{Kategorie:} HK3 -- Soziale Eingebundenheit\\
  \textbf{Begr\"undung:} Die Passage beschreibt eine emotionale Verbindung zur KI, die über rationale Nutzung hinausgeht und auf ein Erleben von Zugehörigkeit und Beziehungsqualität auch gegenüber nicht-menschlichen Akteuren hinweist. Dies berührt die Kategorie der sozialen Eingebundenheit, indem neue – hier quasi-soziale – Beziehungsformen thematisiert werden.
  \\\textit{Memo: Die Analogie zu Star Wars (Mensch-Roboter-Beziehung) ist auffällig: Die Interviewperson positioniert sich selbst als besonders affin für emotionale KI-Beziehungen ('leider für das zu sehr geeignet'), was auf eine individuelle Disposition hinweist, die den motivationalen Umgang mit KI grundlegend färbt. Eine Codierung unter HK0 wäre alternativ denkbar, da es sich auch um eine biografisch-dispositionelle Selbstbeschreibung handelt; HK3 wurde bevorzugt, weil das Beziehungserleben gegenüber KI explizit thematisiert wird.}
\end{description}
\medskip

% ── Unkodierte Passagen ─────────────────────────────────────────
\subsection*{Unkodierte Passagen}

{\small\color{gray}
\noindent Die folgenden Sinneinheiten wurden im ersten Codierdurchgang
keiner Hauptkategorie zugeordnet (kein motivationaler Bezug zum
Untersuchungsgegenstand):
\begin{itemize}
  \item \texttt{\#00:00:10\#} (SE 1): \textit{Bitte starten. Es ist mir eine große Ehre, dass ich dabei sein darf. Geht schon.\ldots} -- Keine kodierrelevante Passage – reine Gesprächseröffnung ohne motivationalen Bezug.
  \item \texttt{\#00:27:53\#} (SE 36): \textit{2010 bin ich in dieser Rolle, seit November 2010.\ldots} -- Keine kodierrelevante Passage. Rein faktische Angabe zur Betriebszugehörigkeit ohne motivationalen Bezug.
  \item \texttt{\#00:28:03\#} (SE 37): \textit{Sechs sind bei mir mit insgesamt 95 FTEs und rund 105 Personen.\ldots} -- Rein faktische Angabe zur Teamgröße (6 Einheiten, 95 FTEs, 105 Personen) ohne motivationalen Bezug.
  \item \texttt{\#00:28:19\#} (SE 38): \textit{Also, das sind die OE-Leiter, dazu noch eine Assistentin, zwei Fachexperten und \ldots} -- Rein faktische Beschreibung der Organisationsstruktur und direkten Berichtslinie ohne motivationalen Bezug.
  \item \texttt{\#00:28:38\#} (SE 39): \textit{Ja, wir sind eine Aktiengesellschaft mit den darunterliegenden 39 Genossenschaft\ldots} -- Rein faktische organisationsstrukturelle Beschreibung ohne motivationalen Bezug.
  \item \texttt{\#00:28:50\#} (SE 40): \textit{Ungefähr 990 sind es in der Landesbank.\ldots} -- Keine kodierrelevante Passage – rein faktische Angabe zur Mitarbeiterzahl ohne motivationalen Bezug.
  \item \texttt{\#00:28:57\#} (SE 41): \textit{Täglich.\ldots} -- Keine kodierrelevante Passage. Einwortige Antwort ohne motivationalen Bezug; vermutlich Frequenzangabe zur KI-Nutzung ohne inhaltliche Ausführung.
\end{itemize}
}
